\section{2차 기계 시스템}

\subsection{매스-스프링-댐퍼 시스템이란 무엇인가?}

매스-스프링-댐퍼(mass-spring-damper) 시스템은 진동과 임피던스를 설명할 때 가장 많이 쓰는 기본 모델이다.
여기서 임피던스는 일단 힘을 주어 움직이게 하려 할 때 시스템이 얼마나, 어떤 방식으로 버티는가를 나타내는 말로 잡아두자.
이름 그대로 세 가지 요소가 함께 움직인다.

\begin{itemize}
  \item 매스(mass): 움직이는 물체이다. 질량 $m$이 클수록 같은 힘으로 속도를 바꾸기 어렵다.
  \item 스프링(spring): 물체를 평형점으로 되돌리려는 요소이다. 강성 $k$가 클수록 같은 변위에서 더 큰 복원력이 생긴다.
  \item 댐퍼(damper): 움직임을 늦추고 진동 에너지를 빼내는 요소이다. 감쇠계수 $b$가 클수록 같은 속도에서 더 큰 저항력이 생긴다.
\end{itemize}

\begin{table}[!htbp]
  \centering
  \caption{매스-스프링-댐퍼 모델에서 먼저 붙잡을 기호}
  \small
  \setlength{\tabcolsep}{4pt}
  \begin{tabularx}{\linewidth}{@{}>{\raggedright\arraybackslash}p{0.16\linewidth}>{\raggedright\arraybackslash}p{0.20\linewidth}>{\raggedright\arraybackslash}X@{}}
    \toprule
    기호 & 단위 & 직관 \\
    \midrule
    $x$ & m & 평형점에서 얼마나 벗어났는지 나타내는 위치 또는 변위 \\
    $\dot{x}$ & m/s & 위치가 얼마나 빠르게 변하는지 나타내는 속도 \\
    $\ddot{x}$ & m/s$^2$ & 속도가 얼마나 빠르게 변하는지 나타내는 가속도 \\
    $m$ & kg & 속도 변화를 싫어하는 정도, 즉 관성의 크기 \\
    $b$ & N\,s/m & 속도에 비례해 운동 에너지를 빼앗는 정도 \\
    $k$ & N/m & 위치가 벗어났을 때 되돌리려는 힘의 기울기 \\
    $f$ & N & 외부에서 물체에 가하는 힘 \\
    \bottomrule
  \end{tabularx}
\end{table}

단위 감각도 잠깐 잡아두자.
$1~\mathrm{N}$은 $1~\mathrm{kg}$ 물체의 속도를 매초 $1~\mathrm{m/s}$씩 바꾸는 힘이다.
$1~\mathrm{N/m}$는 $1~\mathrm{m}$ 변위당 $1~\mathrm{N}$의 힘이 늘어나는 강성이고, $1~\mathrm{N\,s/m}$는 $1~\mathrm{m/s}$ 속도당 $1~\mathrm{N}$의 저항력이 생기는 감쇠계수이다.

방금 본 RLC 회로와 비교하면 이 세 요소의 역할이 더 선명해진다.
전기 회로가 아직 익숙하지 않아도 괜찮다.
지금은 질량은 속도 변화를 싫어하고, 스프링은 위치를 되돌리며, 댐퍼는 운동 에너지를 줄인다는 세 역할만 붙잡으면 된다.
전하 $q$와 전류 $i=\dot{q}$가 전기적 에너지 상태라면, 여기서는 위치 $x$와 속도 $\dot{x}$가 기계적 에너지 상태이다.
전하가 변하는 빠르기가 전류였다면, 위치가 변하는 빠르기는 속도이다.
인덕터 $L$이 전류 변화에 둔감함을 만들었듯이 질량 $m$은 속도 변화에 둔감함을 만든다.
저항 $R$이 전기 에너지를 열로 소산했듯이 댐퍼 $b$는 운동 에너지를 소산한다.
커패시터 항 $\frac{1}{C}q$가 전하를 기준 상태로 되돌리려 했다면, 스프링 항 $kx$는 위치를 평형점으로 되돌리려 한다.

가장 단순한 상황을 생각해보자.
책상 위에 작은 카트가 있고, 카트는 벽에 연결된 스프링과 댐퍼에 붙어 있다.
카트를 오른쪽으로 당겼다가 놓으면 스프링은 카트를 다시 왼쪽으로 잡아당긴다.
하지만 카트는 질량을 가지고 있으므로 평형점에 도착하자마자 멈추지 못하고 지나간다.
그러면 스프링은 반대쪽에서 다시 카트를 잡아당긴다.
이 때문에 카트는 평형점 주변을 왔다 갔다 한다.
이것이 진동이다.

그런데 현실에서는 진동이 영원히 계속되지 않는다.
공기저항, 마찰, 내부 손실, 오일 저항 같은 효과가 에너지를 조금씩 빼앗아 간다.
이 역할을 모델 안에 넣은 것이 댐퍼이다.
댐퍼가 없으면 이상적인 시스템은 계속 흔들리고, 댐퍼가 있으면 진동이 점점 줄어든다.

이 모델을 그릴 때 보통 평형점을 $x=0$으로 둔다.
물체가 오른쪽으로 이동하면 $x>0$, 왼쪽으로 이동하면 $x<0$이다.
속도는 $\dot{x}$, 가속도는 $\ddot{x}$로 쓴다.
외부에서 물체에 가하는 힘은 $f(t)$라고 하자.
그러면 이 시스템의 방정식에는 크게 세 가지 동역학 항이 나타난다.
아래 항들은 잠시 뒤 부호를 정리해서 왼쪽으로 모아 쓸 항들이며, 실제 스프링 힘과 댐퍼 힘의 방향은 바로 다음 문단에서 따로 확인한다.

\begin{align}
  \text{관성항} &\quad m\ddot{x}, \\
  \text{감쇠항} &\quad b\dot{x}, \\
  \text{스프링항} &\quad kx.
\end{align}

이 세 항의 합이 외력과 균형을 이루면 다음 방정식이 나온다.
\begin{equation}
  m\ddot{x}(t)+b\dot{x}(t)+kx(t)=f(t).
  \label{eq:msd_force_balance_intro}
\end{equation}

여기서 왜 스프링 힘과 댐퍼 힘에 음수 부호가 없고 왼쪽에 $+b\dot{x}+kx$로 쓰였는지 헷갈릴 수 있다.
실제로 물체에 작용하는 스프링 힘은 변위를 줄이는 방향이고, 댐퍼 힘은 속도와 반대 방향이므로
\begin{equation}
  f_s=-kx, \qquad f_d=-b\dot{x}
\end{equation}
로 쓴다.
뉴턴의 법칙을 물체에 작용하는 힘의 합으로 쓰면
\begin{equation}
  m\ddot{x}=f+f_s+f_d
  =
  f-kx-b\dot{x}
\end{equation}
이다.
이 식에서 스프링항과 감쇠항을 왼쪽으로 옮기면 식 \eqref{eq:msd_force_balance_intro}처럼
\begin{equation}
  m\ddot{x}+b\dot{x}+kx=f
\end{equation}
가 된다.
이 식은 외워야 할 공식이라기보다, 외력 $f$가 감당해야 할 부담을 나누어 적은 장부에 가깝다.
스프링과 댐퍼의 실제 힘은 여전히 $f_s=-kx$, $f_d=-b\dot{x}$이다.
다만 그 항들을 왼쪽으로 옮겨, 외력이 맞서야 할 항으로 읽는 것이다.
\begin{itemize}
  \item $m\ddot{x}$: 질량을 가속시키는 데 필요한 힘이다. 질량이 클수록 같은 가속도를 만들기 위해 더 큰 힘이 필요하다.
  \item $b\dot{x}$: 움직이는 동안 댐퍼가 요구하는 힘이다. 속도가 클수록 에너지가 더 빨리 소산된다.
  \item $kx$: 평형점에서 벗어난 변위를 유지하거나 더 밀기 위해 필요한 힘이다. 더 멀리 밀수록 스프링 복원력도 커진다.
\end{itemize}

따라서 매스-스프링-댐퍼 시스템은 1자유도 임피던스 직관의 출발점이 된다.
질량, 댐퍼, 스프링은 모두 외부에서 주는 운동에 부담을 주지만, 같은 방식으로 에너지를 없애는 것은 아니다.
질량은 운동 에너지를 저장하고, 스프링은 탄성 퍼텐셜 에너지를 저장하며, 댐퍼만 에너지를 열이나 내부 손실로 소산한다.
즉 관성은 가속도에, 감쇠는 속도에, 강성은 변위에 반응한다.
다음은 뒤에서 다시 만날 기계적 임피던스의 짧은 맛보기다.
운동방정식은 위치 $x$를 중심으로 쓰지만, 기계적 임피던스는 힘을 속도 $\dot{x}$로 나눈다.
위치 $x$는 스프링 에너지를 정하는 저장 변수이고, 속도 $\dot{x}$는 실제로 에너지가 오가는 흐름 변수이기 때문이다.
기계적 임피던스가 크다는 말은 같은 속도로 움직이기 위해 더 큰 힘이 필요하다는 뜻이다.
단위도 힘을 속도로 나눈 것이므로 $\mathrm{N}/(\mathrm{m/s})=\mathrm{N\,s/m}$가 된다.
뒤에서 라플라스 변환을 사용할 때는 외부에서 공급한 힘 $f$를 기준으로 $m\ddot{x}+b\dot{x}+kx=f$를 쓰고, 영 초기조건의 전달함수 문맥에서 속도 $V_x=sX$로 나누어 $ms+b+k/s$를 얻는다.
지금은 외력이 만든 운동을 질량, 댐퍼, 스프링이 서로 다른 방식으로 나누어 감당한다는 점만 가져가도 충분하다.

\subsection{스프링이란 무엇인가?}

스프링(spring)은 늘어나거나 압축되면 원래 길이로 돌아가려는 물체이다.
볼펜 안의 작은 스프링, 자동차 서스펜션, 고무줄, 금속판의 휘어짐은 모두 넓은 의미의 스프링이다.
스프링의 핵심 법칙은 로버트 훅(Robert Hooke)이 17세기에 정리한 훅의 법칙이다.
훅은 1678년에 탄성 물체의 변형과 힘 사이의 비례 관계를 설명했고, 이를 라틴어 ut tensio, sic vis로 표현한 것으로 알려져 있다 \cite{hooke1678depotentiarestitutiva}.
늘어난 만큼 힘도 커진다는 뜻이며, 현대식으로 쓰면 탄성 범위 안에서는 힘이 변위에 비례한다는 말이다.
물론 실제 스프링은 너무 많이 늘리면 영구 변형되거나 끊어진다.
따라서 훅의 법칙은 모든 변형에 대해 무조건 맞는 법칙이 아니라, 재료가 탄성 범위 안에서 비교적 작게 변형될 때 잘 맞는 근사 법칙이다.

스프링을 읽을 때는 먼저 평형점(equilibrium point)을 잡아야 한다.
힘의 방향도, 에너지의 최솟값도 모두 이 기준점에서 결정되기 때문이다.
자연 평형점은 외력이 없을 때 스프링이 자연스럽게 머무는 위치이다.
다만 외력이 일정하게 걸리면 실제로 멈춰 서는 정적 평형점은 이 자연 위치에서 조금 이동할 수 있다.
즉 자연 평형점은 아무 힘도 없을 때의 0점이고, 정적 평형점은 일정한 힘이 계속 걸릴 때 새로 멈추는 위치이다.
예를 들어 오른쪽으로 일정한 힘을 계속 가하면, 스프링이 그 힘과 같은 크기로 왼쪽 복원력을 만들 만큼 늘어난 위치에서 새로 멈춘다.
뒤의 임피던스 제어 장에서 구분하는 목표점 $x_d$와 최종 위치 $x_{\mathrm{ss}}$도 바로 이 차이를 로봇 언어로 옮긴 것이다.

평형점에서 오른쪽으로 $x$만큼 밀었다고 하자.
스프링은 다시 왼쪽으로 잡아당긴다.
반대로 왼쪽으로 $x$만큼 밀면 오른쪽으로 밀어낸다.
즉, 스프링 힘은 항상 변위를 줄이는 방향으로 작용한다.
이 성질을 훅의 법칙이라고 한다.
\begin{equation}
  f_s(t)=-kx(t).
  \label{eq:hooke}
\end{equation}
여기서 $k$는 강성(stiffness)이다.
단위는 N/m이다.
$k=100~\mathrm{N/m}$라는 말은 스프링을 1 m 늘리려면 100 N의 힘이 필요하다는 뜻이다.
물론 실제 작은 스프링을 1 m 늘리는 일은 거의 없으므로, 보통은 1 cm 늘렸을 때 $1~\mathrm{N}$이 필요하다는 식으로 해석하면 된다.
같은 스프링을 2 cm 늘리면 $2~\mathrm{N}$, 5 cm 늘리면 $5~\mathrm{N}$이 필요하다.
이런 선형 비례가 훅의 법칙이 주는 가장 중요한 감각이다.

식 \eqref{eq:hooke}의 음수 부호를 잘 봐야 한다.
$x>0$이면 힘은 음수 방향이고, $x<0$이면 힘은 양수 방향이다.
즉, 스프링 힘은 항상 평형점으로 돌아가려는 복원력(restoring force)이다.
스프링은 변위를 줄이는 복원력으로 나타난다.

스프링을 늘리거나 압축하려면 외부에서 일을 해야 한다.
조금씩 천천히 늘린다고 생각하면, 위치 $x$에서 필요한 힘의 크기는 $kx$이다.
힘-변위 그래프로 보면 더 쉽다.
가로축은 변위 $x$, 세로축은 필요한 힘의 크기 $kx$이다.
스프링은 많이 늘릴수록 필요한 힘이 선형으로 커지므로, 그래프 아래 면적은 밑변이 $x$, 높이가 $kx$인 삼각형이 된다.
따라서 면적은 $\frac{1}{2}x\cdot kx=\frac{1}{2}kx^2$이다.
이 면적이 바로 스프링에 저장되는 일, 즉 에너지이다.
적분에 익숙하지 않다면 이 삼각형 면적만으로 먼저 이해해도 좋다.
적분식은 그 삼각형 면적 계산을 아주 작은 조각들의 합으로 더 엄밀하게 쓴 것이다.
따라서 0에서 $x$까지 늘리는 데 필요한 일은
\begin{equation}
  U_s(x)
  =
  \int_0^x k\xi\,\dd \xi
  =
  \frac{1}{2}kx^2
  \label{eq:spring_energy}
\end{equation}
이다.
이 값이 스프링에 저장된 탄성 퍼텐셜 에너지(elastic potential energy)이다.
왜 $\frac{1}{2}$가 붙는지도 이 적분에서 바로 보인다.
스프링을 처음 늘릴 때는 힘이 거의 필요 없지만, 많이 늘릴수록 더 큰 힘이 필요하다.
그래서 평균적으로 필요한 힘이 최종 힘 $kx$의 절반이 되고, 에너지가 $\frac{1}{2}kx^2$가 된다.

힘과 퍼텐셜 에너지의 관계는 다음과 같다.
\begin{equation}
  f_s(x)=-\frac{\dd U_s(x)}{\dd x}.
\end{equation}
퍼텐셜 에너지가 높은 곳에서 낮은 곳으로 내려가려는 방향이 힘의 방향이다.
음수 부호는 힘이 에너지 언덕을 올라가는 방향이 아니라, 에너지가 낮아지는 방향으로 작용한다는 뜻이다.
스프링의 퍼텐셜 에너지 $U_s=\frac{1}{2}kx^2$는 $x=0$에서 가장 작고, $|x|$가 커질수록 커진다.
그래서 스프링은 항상 $x=0$으로 돌아가려 한다.
이 직관은 나중에 로봇 임피던스 제어로 넘어갈 때 다시 등장한다.
로봇 끝단 주변에 가상의 에너지 지형을 만들고, 그 에너지의 음의 기울기를 힘으로 사용하기 때문이다.

\subsection{댐퍼란 무엇인가?}

댐퍼(damper)는 운동을 천천히 만들고 진동을 줄이는 장치이다.
문이 천천히 닫히게 만드는 도어 클로저, 자동차 쇼크 업소버, 피스톤이 오일 속에서 움직이는 장치가 대표적인 예이다.
스프링이 위치에 반응한다면, 댐퍼는 속도에 반응한다.

역사적으로 댐퍼는 무언가를 움직이게 하는 장치라기보다 너무 심하게 흔들리지 않게 만드는 장치로 발전해 왔다.
마차와 자동차의 서스펜션, 철도 차량, 문 닫힘 장치, 계측기의 바늘 안정화 장치처럼, 실제 기계는 스프링만으로는 충분하지 않았다.
스프링은 충격을 부드럽게 받아낼 수 있지만, 받은 에너지를 다시 되돌려주기 때문에 계속 출렁이게 만든다.
차가 방지턱을 넘을 때 스프링만 있다면 차체는 한참 동안 위아래로 흔들릴 것이다.
쇼크 업소버 같은 유압 댐퍼는 이 운동 에너지를 유체 마찰과 열로 바꾸어 줄여준다.
기계진동 이론에서 점성 감쇠 모델은 이러한 현상을 속도에 비례하는 저항력으로 단순화한 것이다 \cite{libretexts_viscous_damping}.

댐퍼를 가장 쉽게 떠올리는 방법은 주사기나 피스톤을 생각하는 것이다.
공기나 오일이 들어 있는 실린더 안에서 피스톤을 천천히 밀면 비교적 쉽게 움직인다.
하지만 아주 빠르게 밀려고 하면 갑자기 큰 저항이 느껴진다.
왜냐하면 내부의 유체가 좁은 통로나 틈을 지나 이동해야 하기 때문이다.
천천히 움직일 때는 유체가 따라올 시간이 있지만, 빠르게 움직일 때는 유체가 따라오지 못해 큰 압력 차이가 생긴다.
이 압력 차이가 피스톤의 움직임을 방해한다.
이런 장치가 댐퍼의 물리적 이미지이다.

여기서 댐퍼가 스프링과 전혀 다르다는 점을 짚고 넘어가야 한다.
스프링은 위치가 변하면 힘을 만든다.
오른쪽으로 10 cm 밀어 놓고 가만히 잡고 있으면, 스프링은 계속 왼쪽으로 잡아당긴다.
반면 이상적인 댐퍼는 위치만 바뀌어 있고 속도가 0이면 힘을 만들지 않는다.
댐퍼는 어디에 있느냐가 아니라 얼마나 빠르게 움직이느냐에 반응한다.

또한 댐퍼는 단순 마찰과도 조금 다르다.
마른 마찰(Coulomb friction)은 대략 일정한 크기의 힘으로 운동을 방해한다.
속도가 작아도, 커도 마찰력 크기가 거의 일정하다고 보는 모델이다.
반면 점성 감쇠(viscous damping)는 속도가 빠를수록 저항력이 커진다고 보는 모델이다.
유체 속을 움직이는 물체나 오일 댐퍼는 많은 경우 점성 감쇠 모델로 근사할 수 있다.

이상적인 점성 감쇠기(viscous damper)는 다음과 같이 모델링한다.
\begin{equation}
  f_d(t)=-b\dot{x}(t).
  \label{eq:damper_force}
\end{equation}
이 식에서 $b$는 감쇠계수(damping coefficient)이다.
단위는 N s/m이다.
$b=10~\mathrm{N\,s/m}$라는 말은 속도가 $1~\mathrm{m/s}$일 때 $10~\mathrm{N}$의 저항력이 생긴다는 뜻이다.
속도가 $0.1~\mathrm{m/s}$이면 저항력은 $1~\mathrm{N}$이고, 속도가 $2~\mathrm{m/s}$이면 저항력은 $20~\mathrm{N}$이 된다.
가만히 잡고 있어 속도가 0이면, 변위가 아무리 남아 있어도 이상적인 댐퍼 힘은 0이다.
즉, 댐퍼 힘은 속도에 비례한다.

댐퍼 힘도 음수 부호를 가진다.
속도 $\dot{x}>0$이면 물체는 양의 방향으로 움직이고 있으므로, 댐퍼 힘은 음의 방향으로 작용한다.
반대로 $\dot{x}<0$이면 댐퍼 힘은 양의 방향으로 작용한다.
즉, 댐퍼는 항상 운동 방향의 반대 방향으로 힘을 만든다.
댐퍼는 위치가 아니라 속도에 반응해 운동 에너지를 줄인다.

댐퍼에서 꼭 붙잡을 점은 에너지를 저장하지 않는다는 것이다.
스프링은 에너지를 저장했다가 다시 돌려줄 수 있지만, 댐퍼는 운동 에너지를 열, 유체 마찰, 내부 마찰 등으로 바꾸어 없앤다.
여기서 일률(power)은 단위 시간당 에너지가 얼마나 오가는지를 나타내는 양이다.
힘이 물체를 조금 움직이면 일이 생기고, 그 일을 얼마나 빠르게 하고 있는지가 일률이다.
그래서 직선 운동에서는 힘과 속도를 곱하면 단위 시간당 일, 즉 에너지 변화율이 된다.
직선 운동에서는 힘과 속도의 곱
\begin{equation}
  p(t)=f(t)\dot{x}(t)
\end{equation}
로 계산한다.
힘과 속도가 같은 방향이면 물체에 에너지가 들어가고, 반대 방향이면 물체에서 에너지가 빠져나간다.
댐퍼가 물체에 가하는 순간 일률은
\begin{equation}
  p_{\mathrm{on}}(t)=f_d(t)\dot{x}(t)
  =
  -b\dot{x}(t)^2
  \le 0
\end{equation}
이다.
이 값이 항상 0 이하라는 것은 댐퍼가 물체에 에너지를 넣어주는 것이 아니라, 항상 에너지를 빼앗는다는 뜻이다.
반대로 댐퍼가 소산하는 일률을 양수로 쓰면
\begin{equation}
  p_d(t)=b\dot{x}(t)^2\ge 0
\end{equation}
이다.
속도의 제곱이 들어가기 때문에 방향과 상관없이 에너지는 항상 줄어든다.
오른쪽으로 움직여도 에너지가 줄고, 왼쪽으로 움직여도 에너지가 줄어든다.
댐퍼는 운동 방향을 바꾸어 에너지를 되돌려주는 장치가 아니라, 운동 에너지를 열이나 유체 손실로 바꾸어 없애는 장치이다.

그래서 댐퍼는 진동을 가라앉힌다.
진동은 에너지가 계속 왔다 갔다 하면서 생기는 운동인데, 댐퍼가 매 순간 에너지를 조금씩 빼앗아 가면 진폭이 점점 줄어든다.
감쇠가 너무 작으면 진동이 오래 남고, 감쇠가 너무 크면 움직임이 지나치게 둔해진다.
공학에서 좋은 감쇠를 찾는다는 것은 빠르게 움직이면서도 불필요한 진동을 남기지 않는 균형점을 찾는 일이다.

댐퍼의 역할은 로봇 제어에서도 그대로 나타난다.
목표 위치까지 빨리 가고 싶어서 강성을 크게 만들면 로봇은 목표점으로 강하게 끌려간다.
하지만 감쇠가 부족하면 목표점을 지나치고 다시 돌아오며 흔들릴 수 있다.
이때 속도에 비례하는 감쇠항을 추가하면 로봇이 목표점 근처에서 가진 운동 에너지를 줄일 수 있다.
그래서 PD 제어에서 D 항은 단순히 미분항이 아니라, 물리적으로는 가상의 댐퍼 역할을 한다.

\subsection{질량과 관성이란 무엇인가?}

질량(mass)은 물체가 속도 변화를 거부하는 정도이다.
정지한 물체를 움직이게 하거나, 움직이는 물체를 멈추게 하거나, 움직이는 방향을 바꾸려면 힘이 필요하다.
이 성질을 관성(inertia)이라고 부른다.

관성이라는 생각은 뉴턴 역학의 중심에 있다.
아이작 뉴턴(Isaac Newton)은 1687년에 출간한 프린키피아(Philosophiae Naturalis Principia Mathematica)에서 운동 법칙을 정식화했고, 물체가 외력이 없으면 정지 상태나 등속 직선 운동을 유지한다는 관성의 법칙을 제시했다 \cite{newton1687principia}.
일상적으로는 무거운 물체는 움직이기 어렵다고 말하지만, 공학적으로는 질량이 큰 물체는 속도를 바꾸기 어렵다고 말하는 편이 더 정확하다.
정지한 물체를 출발시키는 것도 속도 변화이고, 움직이는 물체를 멈추는 것도 속도 변화이며, 방향을 바꾸는 것도 속도 변화이다.
질량은 바로 이 변화에 대한 저항이다.

뉴턴의 제2법칙은
\begin{equation}
  f_m(t)=m\ddot{x}(t)
\end{equation}
이다.
여기서 $f_m$은 스프링 힘이나 댐퍼 힘처럼 따로 존재하는 새 부품의 힘이라기보다, 질량을 그 가속도로 움직이게 만들기 위해 필요한 순힘을 뜻한다.
같은 힘 $f$를 가했을 때 질량 $m$이 크면 가속도 $\ddot{x}$는 작다.
예를 들어 같은 $10~\mathrm{N}$의 힘을 $1~\mathrm{kg}$ 물체에 가하면 가속도는 $10~\mathrm{m/s^2}$이지만, $10~\mathrm{kg}$ 물체에 가하면 $1~\mathrm{m/s^2}$만 나온다.
따라서 질량은 속도 자체를 방해하는 것이 아니라, 속도가 변하는 것을 방해한다.
이 점에서 질량은 전기 시스템의 인덕터와 닮았다.
인덕터가 전류 변화에 저항하듯, 질량은 속도 변화에 저항한다.

질량에 저장된 운동 에너지는
\begin{equation}
  T(t)=\frac{1}{2}m\dot{x}(t)^2
\end{equation}
이다.
왜 $\frac{1}{2}$가 붙는지도 일의 관점에서 볼 수 있다.
속도를 $v=\dot{x}$라고 쓰면, 질량을 아주 조금 더 빠르게 만드는 데 필요한 미소 일은
\begin{align}
  \dd W
  &=
  f\,\dd x
  =
  m\ddot{x}\,\dd x \nonumber\\
  &=
  m\frac{\dd v}{\dd t}v\,\dd t
  =
  mv\,\dd v.
\end{align}
따라서 정지 상태에서 속도 $v$까지 올리는 동안 저장되는 운동 에너지는
\begin{equation}
  T=\int_0^v mv\,\dd v=\frac{1}{2}mv^2
\end{equation}
가 된다.
속도가 커질수록 같은 속도 증가를 더하는 데 필요한 일이 더 커지므로, 에너지 그래프가 직선이 아니라 포물선이 된다.
속도가 두 배가 되면 운동 에너지는 네 배가 된다.
따라서 빠르게 움직이는 물체를 멈추려면 생각보다 많은 에너지를 제거해야 한다.
이것이 큰 질량과 빠른 속도를 가진 시스템에서 감쇠와 제동이 중요한 이유이다.

전기 회로에서 커패시터가 전기장에 에너지를 저장하고, 인덕터가 전류 변화에 버티며, 저항이 에너지를 열로 소산한다면, 기계 시스템에서는 스프링, 질량, 댐퍼가 같은 세 역할을 나누어 맡는다.
스프링은 위치 변화에 탄성 에너지를 저장하고, 질량은 속도 변화에 관성을 보이며, 댐퍼는 운동 에너지를 열이나 마찰 손실로 빼낸다.
그래서 RLC 회로를 배운 뒤 질량-스프링-댐퍼를 보면 완전히 새 모델을 외우는 것이 아니라, 같은 2차 구조를 눈에 보이는 운동으로 다시 읽는 셈이다.

\begin{table}[!htbp]
  \centering
  \caption{역사적 단서로 기억하는 매스-스프링-댐퍼 요소의 역할}
  \small
  \setlength{\tabcolsep}{3pt}
  \begin{tabularx}{\linewidth}{@{}>{\raggedright\arraybackslash}p{0.13\linewidth}>{\raggedright\arraybackslash}p{0.25\linewidth}>{\raggedright\arraybackslash}p{0.23\linewidth}>{\raggedright\arraybackslash}X@{}}
    \toprule
    항목 & 역사적 또는 공학적 단서 & 물리적 역할 & 임피던스 관점의 기억 문장 \\
    \midrule
    스프링 $k$ & 훅의 법칙: 늘어난 만큼 복원력도 커진다 \cite{hooke1678depotentiarestitutiva}. & 위치 변화에 따라 탄성 에너지를 저장한다. & 평형점에서 멀어질수록 더 큰 힘으로 되돌리려 한다. 전기 쪽 커패시터와 닮지만, 기계 강성은 $k$, 전기 복원 항은 $1/C$로 대응된다. \\
    질량 $m$ & 뉴턴의 운동 법칙: 외력이 없으면 운동 상태를 유지한다 \cite{newton1687principia}. & 속도 변화에 맞서며 운동 에너지를 저장한다. & 속도 자체를 없애려는 요소가 아니라, 속도가 바뀌는 것을 어렵게 만드는 관성 요소이다. 전기 쪽 인덕터와 닮았다. \\
    댐퍼 $b$ & 차량 서스펜션, 도어 클로저, 유체 저항처럼 흔들림을 줄이려는 장치에서 직관을 얻을 수 있다. & 속도에 비례하는 힘으로 운동 에너지를 열이나 유체 손실로 소산한다 \cite{libretexts_viscous_damping}. & 움직일 때만 작동하는 브레이크이다. 빠를수록 더 크게 버티며, 저장한 에너지를 되돌려주지 않는다. \\
    MSD 모델 & 진동을 설명하기 위한 가장 작은 기계 모델이다. & 질량과 스프링이 에너지를 교환하고, 댐퍼가 전체 에너지를 줄인다. & $m$은 지나가게 하고, $k$는 되돌리며, $b$는 남은 에너지를 빼앗는다. 이 세 역할이 모여 고유진동수와 감쇠비를 만든다. \\
    \bottomrule
  \end{tabularx}
  \label{tab:msd-history-role}
\end{table}

\subsection{질량-스프링-댐퍼 모델}

기계 시스템에서 가장 대표적인 2차 시스템은 질량-스프링-댐퍼 모델이다.
여기서 2차(second-order)라는 말은 미분방정식에서 가장 높은 미분이 두 번 미분한 가속도 $\ddot{x}$라는 뜻이다.
또한 미래의 움직임을 예측하려면 현재 위치 $x$만으로는 부족하고, 현재 속도 $\dot{x}$까지 알아야 한다는 뜻이기도 하다.
같은 위치에 있어도 오른쪽으로 빠르게 지나가는 중인지, 왼쪽으로 되돌아오는 중인지에 따라 다음 움직임이 완전히 달라지기 때문이다.
질량 $m$에 외력 $f(t)$가 작용하고, 질량은 스프링 $k$와 감쇠기 $b$에 연결되어 있다고 하자.
뉴턴의 제2법칙을 적용하면
\begin{equation}
  m \ddot{x}(t) + b \dot{x}(t) + k x(t) = f(t)
  \label{eq:msd}
\end{equation}
가 된다.
양변을 $m$으로 나누면
\begin{equation}
  \ddot{x}(t) + \frac{b}{m}\dot{x}(t) + \frac{k}{m}x(t)
  =
  \frac{1}{m}f(t)
\end{equation}
이다.
이제 표준 2차 시스템의 모양과 비교한다.
왼쪽은 시스템 자체의 성질을 담고, 오른쪽은 밖에서 무엇을 넣는지를 담는다.
여기서 표준형으로 바꾸는 이유는 입력보다도 왼쪽 계수에서 고유진동수 $\omega_n$과 감쇠비 $\zeta$를 읽기 위해서이다.
표준형은 보통
\begin{equation}
  \ddot{x}(t)+2\zeta\omega_n\dot{x}(t)+\omega_n^2x(t)
  =
  \frac{1}{m}f(t)
  \label{eq:msd_standard_form}
\end{equation}
처럼 쓴다.
이 식은 힘 입력형으로 쓴 표준형이다.
앞 장처럼 오른쪽을 $\omega_n^2 u$ 모양으로 맞추고 싶다면 $u=f/k$로 두면 된다.
이때 $u$는 그 힘이 정적으로 만들 변위, 즉 스프링만 놓고 오래 기다렸을 때의 위치라고 보면 된다.
앞 장에서 $u=Cv_{\mathrm{in}}$로 전압 입력을 전하 단위의 기준 입력으로 바꾸어 읽었듯, 여기서는 $u=f/k$로 힘 입력을 오래 기다렸을 때의 정적 변위 단위로 바꾸어 읽는다.
이렇게 쓰는 이유는 $\omega_n$과 $\zeta$가 응답의 모양을 바로 말해주는 두 숫자이기 때문이다.
식의 계수를 서로 맞추면
\begin{align}
  \omega_n &= \sqrt{\frac{k}{m}}, \\
  \zeta &= \frac{b}{2\sqrt{mk}}.
\end{align}
즉, $\omega_n^2$은 위치를 되돌리는 항의 세기이고, $2\zeta\omega_n$은 속도에 비례해 진동을 줄이는 항의 세기이다.

이 식은 1자유도 임피던스 제어 직관의 출발점이 된다.
강성 $k$가 크면 위치 오차에 대해 큰 힘이 발생한다.
감쇠 $b$가 크면 속도에 비례한 저항이 커진다.
질량 $m$이 크면 같은 힘에 대해 가속도가 작아지고, 시스템은 더 둔하게 움직인다.

\subsection{진동의 원리: 에너지 교환}

질량-스프링-댐퍼 시스템의 진동은 운동 에너지와 탄성 퍼텐셜 에너지 사이의 교환으로 이해할 수 있다.
이제 진동을 힘의 관점이 아니라 에너지 장부로 다시 써 보자.
어디에 에너지가 저장되고, 어디에서 빠져나가는지만 보면 진동이 왜 줄어드는지 훨씬 선명해진다.
총 기계 에너지를
\begin{equation}
  E(t)
  =
  \frac{1}{2}m\dot{x}(t)^2
  +
  \frac{1}{2}kx(t)^2
\end{equation}
라고 하자.
여기서 하고 싶은 일은 에너지의 정확한 값 자체보다, 시간이 지나며 에너지가 늘어나는지 줄어드는지를 보는 것이다.
그래서 에너지의 시간 변화율인 $\dot{E}$를 계산한다.
시간 미분을 계산하면
\begin{align}
  \dot{E}(t)
  &=
  m\dot{x}(t)\ddot{x}(t)
  +
  kx(t)\dot{x}(t) \\
  &=
  \dot{x}(t)\left(m\ddot{x}(t)+kx(t)\right).
\end{align}
식 \eqref{eq:msd}에서
\begin{equation}
  m\ddot{x}(t)+kx(t)=f(t)-b\dot{x}(t)
\end{equation}
이므로
\begin{equation}
  \dot{E}(t)
  =
  f(t)\dot{x}(t)
  -
  b\dot{x}(t)^2.
\end{equation}
첫 번째 항은 외력이 시스템에 넣어주는 일률이고, 두 번째 항은 감쇠기가 소산하는 일률이다.
여기서 스프링과 질량 사이의 에너지 교환은 총에너지 안에서 서로 주고받는 내부 이동이다.
총량을 바꾸는 것은 바깥에서 넣는 일률 $f\dot{x}$와 댐퍼가 빼앗는 일률 $b\dot{x}^2$뿐이다.
외력이 없고 감쇠도 없다면 $\dot{E}=0$이므로 총 에너지는 보존된다.
이때 에너지는 스프링의 퍼텐셜 에너지와 질량의 운동 에너지 사이를 계속 오가며 진동을 만든다.
외력이 없고 감쇠가 있으면 매 순간 $b\dot{x}^2$의 일률로 에너지가 소산되므로 진동 진폭이 점차 줄어든다.
반대로 외부에서 주기적으로 힘을 넣고 있다면 이야기가 달라질 수 있다.
입력 일률 $f\dot{x}$가 감쇠 소산 $b\dot{x}^2$보다 크면, 댐퍼가 에너지를 빼앗고 있어도 전체 에너지는 늘어날 수 있다.
공진에서 큰 응답이 생기는 이유도 외부 입력이 내부 에너지 교환 리듬과 맞아 에너지를 계속 보충하기 때문이다.

한 주기 동안의 에너지 흐름을 천천히 따라가 보자.
카트가 오른쪽 끝까지 밀려 잠깐 멈춘 순간에는 속도가 거의 0이므로 운동 에너지는 작고, 스프링이 많이 늘어나 있어 탄성 퍼텐셜 에너지가 크다.
카트가 평형점을 지날 때는 스프링 변위가 0에 가까워 퍼텐셜 에너지는 작고, 속도가 가장 커져 운동 에너지가 크다.
그 뒤 카트가 왼쪽 끝에 도달하면 다시 속도는 0에 가까워지고, 스프링은 반대 방향으로 변형되어 퍼텐셜 에너지를 저장한다.
댐퍼가 없다면 이 교환이 계속되고, 댐퍼가 있으면 매번 조금씩 에너지가 열로 빠져나가 왕복 폭이 줄어든다.
그림~\ref{fig:msd-energy-exchange}는 이 과정을 카트-스프링-댐퍼의 운동과 에너지 그래프로 함께 보여준다.

\begin{figure}[!htbp]
  \centering
  \resizebox{\linewidth}{!}{%
  \begin{tikzpicture}[
      snap/.style={draw=black!35, rounded corners=2pt, fill=black!2, minimum width=4.25cm, minimum height=2.55cm},
      cart/.style={draw=black!70, fill=blue!8, minimum width=0.85cm, minimum height=0.55cm},
      label/.style={font=\scriptsize, align=center},
      axis/.style={-{Latex[length=2mm]}, draw=black!65},
      curve/.style={thick},
      wall/.style={draw=black!70, line width=0.5pt}
    ]
    \node[snap] (a) at (0,0) {};
    \node[snap] (b) at (4.8,0) {};
    \node[snap] (c) at (9.6,0) {};

    \foreach \x/\title in {0/{최대 변위},4.8/{평형점 통과},9.6/{감쇠 후}} {
      \node[label, font=\small\bfseries] at (\x,1.05) {\title};
      \draw[wall] (\x-1.75,-0.65) -- (\x-1.75,0.55);
      \draw[black!40] (\x-1.95,-0.65) -- (\x-1.55,-0.65);
      \draw[black!40] (\x-1.95,-0.45) -- (\x-1.55,-0.45);
      \draw[black!40] (\x-1.95,-0.25) -- (\x-1.55,-0.25);
      \draw[black!45] (\x-1.8,-0.8) -- (\x+1.85,-0.8);
    }

    \node[cart] (cartA) at (0.85,-0.35) {$m$};

    \draw[black!70] (-1.75,-0.1)
      -- (-1.35,-0.1) -- (-1.2,0.12) -- (-0.95,-0.32) -- (-0.7,0.12)
      -- (-0.45,-0.32) -- (-0.2,0.12) -- (0.05,-0.32) -- (0.3,0.12)
      -- (0.55,-0.1) -- (0.42,-0.1);
    \draw[black!60] (-1.75,-0.52) -- (-1.25,-0.52) -- (-1.25,-0.28) -- (-0.95,-0.28) -- (-0.95,-0.76) -- (-0.65,-0.76) -- (-0.65,-0.52) -- (0.42,-0.52);
    \node[label] at (0,-1.14) {스프링 에너지 큼\\운동 에너지 작음};

    \node[cart] (cartB) at (4.8,-0.35) {$m$};
    \draw[black!70] (3.05,-0.1)
      -- (3.38,-0.1) -- (3.5,0.08) -- (3.68,-0.28) -- (3.86,0.08)
      -- (4.04,-0.28) -- (4.22,0.08) -- (4.4,-0.28) -- (4.58,-0.1);
    \draw[black!60] (3.05,-0.52) -- (3.45,-0.52) -- (3.45,-0.28) -- (3.72,-0.28) -- (3.72,-0.76) -- (3.99,-0.76) -- (3.99,-0.52) -- (4.38,-0.52);
    \draw[-{Latex[length=2mm]}, thick, red!65!black] (5.35,-0.35) -- (6.15,-0.35) node[right, label] {$\dot{x}$ 큼};
    \node[label] at (4.8,-1.14) {스프링 에너지 작음\\운동 에너지 큼};

    \node[cart] (cartC) at (9.78,-0.35) {$m$};
    \draw[black!70] (7.85,-0.1)
      -- (8.15,-0.1) -- (8.28,0.05) -- (8.45,-0.24) -- (8.62,0.05)
      -- (8.79,-0.24) -- (8.96,0.05) -- (9.13,-0.24) -- (9.3,-0.1);
    \draw[black!60] (7.85,-0.52) -- (8.23,-0.52) -- (8.23,-0.28) -- (8.48,-0.28) -- (8.48,-0.76) -- (8.73,-0.76) -- (8.73,-0.52) -- (9.3,-0.52);
    \node[label] at (9.6,-1.17) {감쇠가 에너지를\\열과 손실로 소산};

    \begin{scope}[shift={(0,-4.85)}]
      \draw[axis] (0,0) -- (10.2,0) node[right, label] {시간 $t$};
      \draw[axis] (0,0) -- (0,2.75) node[above, label] {에너지};
      \draw[curve, blue!70!black] plot[smooth] coordinates {
        (0.15,2.15) (0.8,0.55) (1.45,1.75) (2.1,0.45) (2.75,1.35)
        (3.4,0.35) (4.05,1.0) (4.7,0.28) (5.35,0.72) (6.0,0.22)
        (6.65,0.52) (7.3,0.17) (8.0,0.36) (8.8,0.13) (9.7,0.22)
      };
      \draw[curve, red!70!black, dashed] plot[smooth] coordinates {
        (0.15,0.2) (0.8,1.75) (1.45,0.45) (2.1,1.35) (2.75,0.32)
        (3.4,1.0) (4.05,0.25) (4.7,0.72) (5.35,0.18) (6.0,0.52)
        (6.65,0.13) (7.3,0.36) (8.0,0.1) (8.8,0.22) (9.7,0.08)
      };
      \draw[curve, black!70] plot[smooth] coordinates {
        (0.15,2.35) (1.0,2.18) (2.0,1.85) (3.0,1.5) (4.0,1.18)
        (5.0,0.92) (6.0,0.7) (7.0,0.52) (8.0,0.38) (9.7,0.25)
      };
      \node[label, text=blue!70!black] at (1.25,2.45) {$U_s(t)$: 스프링 에너지};
      \node[label, text=red!70!black] at (6.9,1.15) {$T(t)$: 운동 에너지};
      \node[label, text=black!70] at (8.15,0.75) {$E(t)=T+U_s$};
      \node[label, align=left] at (4.0,2.45) {감쇠가 있으면\\총 에너지가 감소};
    \end{scope}
  \end{tikzpicture}%
  }
  \caption{질량-스프링-댐퍼 시스템의 자유진동은 운동 에너지 $T(t)=\frac{1}{2}m\dot{x}^2$와 탄성 퍼텐셜 에너지 $U_s(t)=\frac{1}{2}kx^2$가 서로 교환되는 과정으로 볼 수 있다. 외력이 없는 자유응답에서 댐퍼가 있으면 $b\dot{x}^2$의 일률로 에너지가 소산되므로 총 기계 에너지 $E(t)$가 줄어들고, 그 결과 진동 진폭도 점차 작아진다.}
  \label{fig:msd-energy-exchange}
\end{figure}


\subsection{고유진동수란 무엇인가?}

여기부터는 입력을 계속 넣었을 때의 응답이 아니라, 처음에 한 번 당기거나 밀어 놓은 뒤 시스템이 스스로 어떻게 움직이는지를 본다.
이렇게 관점을 나누어 두면 전달함수로 보는 강제응답과, 에너지가 오가며 생기는 자유진동이 섞이지 않는다.

고유진동수(natural frequency)는 시스템이 스스로 흔들리고 싶어 하는 속도이다.
여기서 속도는 물체가 앞으로 몇 $\mathrm{m/s}$로 움직인다는 뜻이 아니라, 1초에 몇 번쯤 왕복하려는지에 가까운 흔들림의 빠르기이다.
질량-스프링 시스템을 생각해보자.
물체를 오른쪽으로 당겼다가 손을 놓으면 스프링은 물체를 왼쪽으로 끌어당긴다.
물체가 평형점에 도착했을 때 스프링 힘은 0이지만, 물체는 이미 속도를 가지고 있으므로 그대로 지나간다.
그러면 반대쪽에서 스프링이 다시 물체를 끌어당긴다.
이 과정이 반복되면 진동이 생긴다.

여기서 고유라는 말이 중요하다.
외부에서 계속 흔들어 주지 않아도, 시스템 내부의 질량과 강성만으로 정해지는 자연스러운 흔들림 속도라는 뜻이다.
기타 줄을 튕기면 줄은 자기 고유한 음으로 떨린다.
건물도, 다리도, 자동차 서스펜션도, 로봇 팔도 각각 고유한 흔들림 주파수를 가진다.
이 주파수는 물체가 무거운지, 얼마나 단단하게 지지되어 있는지에 따라 달라진다.
책상 끝에 자를 걸쳐 놓고 튕기는 예도 비슷하다.
자를 짧게 내밀면 더 단단하게 지지된 것처럼 빠르게 떨고, 길게 내밀면 더 부드럽게 지지된 것처럼 느리게 떤다.

감쇠와 외력이 없는 경우 방정식은
\begin{equation}
  m\ddot{x}+kx=0
\end{equation}
이다.
이 식을
\begin{equation}
  \ddot{x}+\frac{k}{m}x=0
\end{equation}
로 나누어 쓰면, 고유진동수는
\begin{equation}
  \omega_n=\sqrt{\frac{k}{m}}
\end{equation}
이다.
단위는 rad/s이다.
왜 제곱근이 나오는지도 간단히 볼 수 있다.
감쇠가 없는 왕복 운동은 일정한 모양으로 반복되므로, 가장 기본적인 반복 함수인 사인이나 코사인으로 시험해 볼 수 있다.
감쇠가 없는 진동이 $x(t)=A\cos(\omega t)$처럼 생겼다고 하면, 두 번 미분한 가속도는 $\ddot{x}(t)=-\omega^2x(t)$가 된다.
이를 $m\ddot{x}+kx=0$에 넣으면
\begin{equation}
  -m\omega^2x+kx=0
\end{equation}
이므로, $x\neq0$인 진동에서는 $\omega^2=k/m$이어야 한다.
따라서 $\omega=\sqrt{k/m}$가 된다.
단위로 보아도 $k/m$은
\begin{equation}
  \frac{\mathrm{N/m}}{\mathrm{kg}}
  =
  \frac{\mathrm{kg\,m/s^2}/\mathrm{m}}{\mathrm{kg}}
  =
  \frac{1}{\mathrm{s^2}}
\end{equation}
이므로, 제곱근을 취하면 $1/\mathrm{s}$가 된다.
rad는 각도를 세는 무차원 단위처럼 취급되므로 이것을 rad/s라고 부른다.
한 주기, 즉 한 바퀴의 위상 변화가 $2\pi$ rad이므로 rad/s를 $2\pi$로 나누면 초당 몇 번 반복되는지 알 수 있다.
일상적인 Hz 단위 주파수 $f_n$으로 바꾸면
\begin{equation}
  f_n=\frac{\omega_n}{2\pi}
\end{equation}
이다.

이 식은 물리적으로 매우 직관적이다.
스프링이 $x$만큼 밀렸을 때 만드는 힘의 크기는 대략 $kx$이고, 이 힘이 질량 $m$을 되돌리는 가속도는 $a\approx(k/m)x$이다.
따라서 $k/m$은 같은 변위가 생겼을 때 얼마나 큰 가속도로 되돌리려 하는지를 나타낸다.
\begin{itemize}
  \item 강성 $k$가 커지면 스프링이 더 세게 잡아당기므로 더 빠르게 흔들린다. 따라서 $\omega_n$은 커진다.
  \item 질량 $m$이 커지면 같은 힘에 대해 가속도가 작아지므로 더 느리게 흔들린다. 따라서 $\omega_n$은 작아진다.
\end{itemize}

고유진동수는 시스템 내부의 복원하려는 성질과 움직임을 바꾸기 싫어하는 성질의 비율에서 나온다.
스프링은 되돌리려 하고, 질량은 속도 변화를 거부한다.
이 둘이 서로 밀고 당기며 만들어내는 고유한 리듬이 고유진동수이다.

숫자로 보면 더 분명하다.
질량이 $m=1~\mathrm{kg}$이고 스프링 강성이 $k=100~\mathrm{N/m}$라면
\begin{equation}
  \omega_n=\sqrt{\frac{100}{1}}=10~\mathrm{rad/s}
\end{equation}
이다.
이를 Hz로 바꾸면
\begin{equation}
  f_n=\frac{10}{2\pi}\approx 1.59~\mathrm{Hz}
\end{equation}
이므로 1초에 약 1.59번 흔들리는 리듬을 가진다.
한 주기로 보면
\begin{equation}
  T=\frac{2\pi}{\omega_n}\approx0.63~\mathrm{s}
\end{equation}
이다.
즉, 이 시스템은 감쇠가 거의 없다면 약 0.63초마다 한 번 왕복하는 리듬을 가진다.
같은 스프링에 질량을 $4~\mathrm{kg}$으로 늘리면
\begin{equation}
  \omega_n=\sqrt{\frac{100}{4}}=5~\mathrm{rad/s}
\end{equation}
가 된다.
질량을 네 배로 늘렸더니 고유진동수는 절반이 되었다.
반대로 질량은 그대로 두고 스프링 강성을 네 배로 키우면 고유진동수는 두 배가 된다.
이처럼 고유진동수는 $k$와 $m$에 선형으로 비례하지 않고 제곱근으로 변한다.
한 줄로 요약하면, 강성 4배는 고유진동수 2배, 질량 4배는 고유진동수 절반이다.

공학적으로 고유진동수는 이 시스템을 어느 속도로 흔들면 위험하거나 크게 반응할 수 있는가를 알려준다.
외부에서 들어오는 반복 힘의 주파수가 고유진동수 근처에 있으면 공진(resonance)이 생길 수 있다.
감쇠가 작고 입력이 충분히 오래 지속되면, 작은 힘을 반복해서 주어도 에너지가 같은 리듬으로 쌓여 응답이 크게 커질 수 있다.
그네를 밀 때를 떠올리면 쉽다.
그네가 돌아오는 타이밍에 맞추어 조금씩 밀면 작은 힘도 계속 에너지를 더해 진폭을 키운다.
반대로 타이밍이 어긋나면 힘을 써도 에너지가 잘 쌓이지 않는다.
그래서 다리, 건물, 자동차, 로봇, 전기회로 모두 고유진동수를 고려해 설계한다.
로봇 제어에서도 강성을 크게 잡으면 목표점으로 빨리 돌아가려는 힘은 커지지만, 동시에 고유진동수도 올라간다.
감쇠가 충분하지 않으면 시스템은 더 빠르게, 더 날카롭게 흔들릴 수 있다.
따라서 강성을 키우면 무조건 좋다는 식으로 볼 수 없고, 강성과 감쇠를 함께 봐야 한다.
다만 여기서의 $\omega_n$은 감쇠가 없다고 보았을 때의 고유진동수이다.
감쇠가 있는 실제 강제응답에서 가장 크게 흔들리는 공진 피크는 보통 이 값 근처에 있지만, 정확히 같은 주파수라고 단정하면 안 된다.
감쇠가 작을 때는 피크가 뚜렷하고 $\omega_n$ 근처에 나타나지만, 감쇠가 커지면 피크가 낮아지고 더 낮은 주파수 쪽으로 이동하거나 아예 뚜렷한 공진 피크가 사라질 수 있다.
처음에는 고유진동수 근처에서 큰 응답이 나올 수 있다고 잡아두자.
정확한 피크 위치와 크기는 감쇠까지 함께 보아야 한다.
심화로 덧붙이면, 힘-변위 전달함수 기준에서 감쇠가 작을 때 변위 응답의 피크는 $\omega_n$보다 약간 낮은 $\omega_n\sqrt{1-2\zeta^2}$ 근처에 생기고, $\zeta\ge1/\sqrt{2}$에서는 뚜렷한 변위 공진 피크가 사라진다.

고유진동수를 직관적으로 요약하면 다음과 같다.
강성은 되돌리는 힘의 세기이고, 질량은 그 힘에 대한 둔감함이다.
강성이 질량에 비해 크면 시스템은 짧은 시간 안에 방향을 바꿀 수 있으므로 빠르게 흔들린다.
질량이 강성에 비해 크면 같은 복원력으로 속도를 바꾸기 어려워 천천히 흔들린다.
그래서 $\omega_n=\sqrt{k/m}$은 단순한 암기식이 아니라, 되돌리는 성질과 버티는 성질의 비율이다.

하지만 실제 시스템에서는 이 리듬이 영원히 계속되지 않는다.
마찰, 댐퍼, 제어 입력이 에너지를 빼앗기 때문이다.
그래서 다음으로 볼 값은 얼마나 빠르게 흔들리는가가 아니라, 그 흔들림이 얼마나 빨리 사라지는가이다.

\subsection{감쇠비란 무엇인가?}

감쇠비(damping ratio) $\zeta$는 현재 시스템 기준에서 감쇠가 얼마나 충분한지를 무차원으로 나타낸 값이다.
감쇠계수 $b$ 자체도 감쇠의 크기를 말해주지만, $b$만 보고는 시스템이 많이 감쇠되는지 적게 감쇠되는지 판단하기 어렵다.
$b=10~\mathrm{N\,s/m}$라는 말은 물체가 $1~\mathrm{m/s}$로 움직일 때 댐퍼가 대략 $10~\mathrm{N}$의 반대 힘을 만든다는 뜻이다.
질량과 스프링 강성이 다르면 같은 $b$라도 효과가 달라지기 때문이다.

예를 들어 아주 가벼운 장난감 카트에 $b=10~\mathrm{N\,s/m}$인 댐퍼를 달면 매우 강한 감쇠처럼 느껴질 수 있다.
하지만 무거운 자동차 차체에 같은 댐퍼를 달면 거의 약한 감쇠처럼 보일 수 있다.
같은 댐퍼라도 시스템이 얼마나 무거운지, 스프링이 얼마나 강한지에 따라 효과가 달라진다.
그래서 감쇠계수 $b$를 그대로 비교하지 않고, 그 시스템에 맞는 기준값으로 나누어 비교한다.
그 기준화된 값이 감쇠비이다.

이때 비교 기준으로 쓰는 값이 임계감쇠(critical damping)이다.
임계감쇠는 복소근이 만들어 내는 반복 진동이 생기는 응답과, 반복 진동 없이 돌아오는 응답 사이의 경계라고 볼 수 있다.
질량-스프링-댐퍼 시스템에서 임계감쇠계수는
\begin{equation}
  b_c=2\sqrt{mk}
\end{equation}
이다.
먼저 말로는 흔들림이 막 사라지는 경계라고 이해하면 된다.
수학적으로는 뒤에서 자세히 볼 특성근과 판별식으로 이 경계를 확인할 수 있다.
감쇠가 있는 자유응답의 특성방정식은
\begin{equation}
  ms^2+bs+k=0
\end{equation}
이고, 이 방정식의 판별식은 $b^2-4mk$이다.
두 근이 복소수이면 흔들리고, 서로 다른 실수이면 흔들리지 않는다.
그 경계는 판별식이 0일 때이므로
\begin{equation}
  b^2-4mk=0
\end{equation}
이다.
$b^2=4mk$에서 제곱근을 취하면 $b=\pm2\sqrt{mk}$인데, 감쇠계수는 물리적으로 양수이므로 양수 근을 택해 $b_c=2\sqrt{mk}$가 된다.
\vmark{low-bc-positive-root}%
감쇠비는 실제 감쇠계수 $b$를 임계감쇠계수 $b_c$로 나눈 값이다.
\begin{equation}
  \zeta=\frac{b}{b_c}
  =
  \frac{b}{2\sqrt{mk}}.
\end{equation}
감쇠비는 임계감쇠에 비해 현재 감쇠가 몇 배인지를 나타낸다.
$b_c$의 단위는 $\mathrm{N\,s/m}$이고, 따라서 $b/b_c$인 $\zeta$는 단위가 없는 숫자이다.

임계감쇠가 왜 기준이 되는지도 중요하다.
대표적인 경우, 감쇠가 너무 작으면 물체는 평형점을 지나쳐 반대쪽으로 넘어간다.
이것이 오버슈트이다.
이후 다시 돌아오고, 또 지나치고, 이런 식으로 진동이 남는다.
감쇠를 조금씩 키우면 이 진동이 줄어든다.
어느 순간부터는 복소근에 의한 반복 진동이 사라진다.
바로 그 경계가 임계감쇠이다.
따라서 표준 1자유도 선형 자유응답에서 임계감쇠는 진동을 없애기 위한 최소한의 경계 감쇠로 이해할 수 있다.
여기서 자유응답은 한 번 당겨 놓고 손을 뗀 뒤 외부 힘 없이 움직이는 응답을 말한다.
계단응답은 목표값이나 입력을 갑자기 바꾸고 그 값을 계속 유지했을 때의 응답을 말한다.
영 초기조건은 시작할 때 위치와 속도 같은 초기 상태를 0으로 두는 계산 조건이다.
대표적인 영 초기조건 계단응답이나 정지 상태에서 변위만 주고 놓는 자유응답에서는, 임계감쇠가 진동 없이 빨리 돌아오는 경계가 된다.
다만 처음 속도가 크면 임계감쇠라도 한 번 평형점을 지날 수 있다.
중요한 차이는 반복해서 왕복하는 진동이 남느냐이고, 초기조건이나 최적화 기준이 달라지면 가장 좋은 감쇠값도 달라질 수 있다.

\begin{table}[!htbp]
  \centering
  \caption{응답 용어를 처음 읽을 때의 짧은 해석}
  \small
  \setlength{\tabcolsep}{4pt}
  \begin{tabularx}{\linewidth}{@{}>{\raggedright\arraybackslash}p{0.22\linewidth}>{\raggedright\arraybackslash}X@{}}
    \toprule
    용어 & 처음 붙잡을 해석 \\
    \midrule
    자유응답 & 한 번 당기거나 밀어 놓고 손을 뗀 뒤, 외부 입력 없이 스스로 움직이는 응답 \\
    계단응답 & 목표값이나 입력을 갑자기 바꾸고, 그 값을 계속 유지했을 때의 응답 \\
    영 초기조건 & 시작 전 위치와 속도 같은 초기 상태를 0으로 두는 계산 조건 \\
    임계감쇠 & 대표 조건에서 반복 진동이 막 사라지는 경계 감쇠 \\
    \bottomrule
  \end{tabularx}
\end{table}

\begin{itemize}
  \item $\zeta=0$이면 감쇠가 전혀 없다. 에너지가 빠져나가지 않으므로 진동이 계속된다.
  \item $0<\zeta<1$이면 감쇠가 있긴 하지만 충분히 크지는 않다. 그래서 진동하면서 점점 줄어든다.
  \item $\zeta=1$이면 임계감쇠이다. 대표적인 조건에서 진동 없이 빠르게 평형점으로 돌아오는 경계이다.
  \item $\zeta>1$이면 과감쇠이다. 진동은 없지만 대표적인 위치 응답에서는 임계감쇠보다 느려질 수 있으며, 안전이나 오버슈트 억제가 더 중요할 때는 의도적으로 선택하기도 한다.
\end{itemize}

자동차 서스펜션을 생각하면 감쇠비의 의미가 잘 보인다.
감쇠가 너무 작으면 차가 방지턱을 넘은 뒤에도 계속 출렁인다.
감쇠가 너무 크면 출렁이진 않지만 충격을 부드럽게 흡수하지 못하고 움직임이 둔해진다.
좋은 설계는 보통 너무 많이 흔들리지 않으면서도 너무 느리지 않은 중간 영역을 찾는 것이다.

조금 더 깊게 보면, 감쇠비는 에너지가 한 번 흔들릴 때마다 얼마나 빠져나가는지와도 연결된다.
다만 이 부분은 그래프의 포락선과 함께 보면 훨씬 덜 추상적이다.
진동하는 부족감쇠 시스템에서는 스프링의 퍼텐셜 에너지와 질량의 운동 에너지가 계속 오가고, 댐퍼가 매 순간 $b\dot{x}^2$의 일률로 에너지를 소산한다.
감쇠가 작으면 한 주기 동안 사라지는 에너지가 적어서 오래 흔들린다.
다만 에너지가 초 단위 시간에서 얼마나 빨리 줄어드는지는 감쇠비 $\zeta$만이 아니라 포락선 감쇠율 $\zeta\omega_n$에도 달려 있다.
부족감쇠 자유응답에서 진폭 포락선은 $e^{-\zeta\omega_n t}$, 에너지 포락선은 대략 $e^{-2\zeta\omega_n t}$로 줄어든다.
따라서 감쇠비는 한 고유 리듬에 대한 상대적 감쇠 정도를 말해주는 지표이고, 실제 정착 속도는 $\omega_n$과 함께 보아야 한다.
$\zeta\ge1$에서는 진동 주기 자체가 사라지므로, 한 주기당 에너지 손실보다 뒤에서 볼 실수 특성근이 만드는 느린 수렴 속도로 해석해야 한다.

감쇠비는 속도 자체를 줄이는 손잡이가 아니다.
현재 시스템의 진동성에 비해 감쇠가 충분한지를 나타내는 상대 지표에 가깝다.
같은 댐퍼라도 무거운 물체와 약한 스프링에 붙으면 응답이 다르게 보이고, 가벼운 물체와 강한 스프링에 붙으면 전혀 다른 응답이 된다.
그래서 공학자는 감쇠계수 $b$만 보지 않고, $m$과 $k$를 함께 묶은 $\zeta=b/(2\sqrt{mk})$를 본다.
감쇠비는 댐퍼 하나의 성질이 아니라 시스템 전체의 성질이다.

로봇 제어에서 감쇠비는 조작감과 안정성에 바로 연결된다.
강성이 큰데 감쇠비가 낮으면 로봇 끝단은 목표점으로 강하게 끌려가지만 목표점을 지나쳐 흔들릴 수 있다.
감쇠비가 너무 높으면 흔들림은 적지만 움직임이 둔하고 무겁게 느껴진다.
그래서 임피던스 제어에서는 원하는 강성 $k$만 정하는 것으로 충분하지 않다.
1축 모델이나 로봇 끝단의 한 방향을 스칼라 유효질량으로 근사할 수 있다면, 원하는 감쇠비 $\zeta$를 정하고
\begin{equation}
  b=2\zeta\sqrt{mk}
\end{equation}
로 감쇠계수까지 함께 정해야 한다.
여기서 로봇 문맥의 $m$은 실제 전체 로봇 질량이 아니라, 그 방향에서 느껴지는 유효질량 또는 교육용 명목 질량으로 읽어야 한다.
예를 들어 $m=1~\mathrm{kg}$, $k=100~\mathrm{N/m}$에서 임계감쇠 $\zeta=1$을 원하면
\begin{equation}
  b=2\cdot 1\cdot \sqrt{1\cdot 100}=20~\mathrm{N\,s/m}
\end{equation}
이다.
조금 더 빠르고 약간 출렁이는 응답을 허용하여 $\zeta=0.7$을 원하면
\begin{equation}
  b=2\cdot 0.7\cdot 10=14~\mathrm{N\,s/m}
\end{equation}
로 잡을 수 있다.
반대로 같은 $b=10~\mathrm{N\,s/m}$이라도 $m=1~\mathrm{kg}$, $k=100~\mathrm{N/m}$이면
\begin{equation}
  \zeta=\frac{10}{2\sqrt{1\cdot100}}=0.5
\end{equation}
이다.
하지만 $m=100~\mathrm{kg}$, $k=100~\mathrm{N/m}$이면
\begin{equation}
  \zeta=\frac{10}{2\sqrt{100\cdot100}}=0.05
\end{equation}
가 된다.
같은 댐퍼인데도 첫 번째 시스템에서는 꽤 의미 있는 감쇠이고, 두 번째 시스템에서는 매우 작은 감쇠처럼 보인다.
이 예시가 감쇠비가 댐퍼 하나의 성질이 아니라 시스템 전체의 성질이라는 점을 잘 보여준다.

강성만 바꾸어도 같은 일이 생긴다.
$m=1~\mathrm{kg}$, $k=100~\mathrm{N/m}$, $b=20~\mathrm{N\,s/m}$이면 $\zeta=1$이다.
그런데 질량과 댐퍼는 그대로 두고 강성만 $k=400~\mathrm{N/m}$로 네 배 키우면
\begin{equation}
  \zeta=\frac{20}{2\sqrt{1\cdot400}}=0.5
\end{equation}
가 된다.
강성이 커지면서 고유진동수는 두 배 빨라졌지만, 같은 댐퍼가 상대적으로는 절반만큼의 감쇠비를 제공하게 된 것이다.
같은 감쇠비 $\zeta=1$을 유지하려면
\begin{equation}
  b=2\sqrt{1\cdot400}=40~\mathrm{N\,s/m}
\end{equation}
처럼 감쇠계수도 두 배로 키워야 한다.
그래서 제어 튜닝에서는 강성을 올릴 때 감쇠도 함께 다시 계산하는 습관을 들여야 한다.
강성 $k$를 올리면 같은 감쇠비를 유지하기 위해 $b$도 $\sqrt{k}$에 비례해 다시 올려야 한다는 뜻이다.
일단은 로봇 끝단의 한 방향만 떼어 1자유도 질량-스프링-댐퍼처럼 본다고 생각하면 된다.
이 장의 식들은 $m>0$, $k>0$, $b\ge0$인 1자유도 선형 점성감쇠 모델, 또는 모드별로 분리해서 볼 수 있는 좌표에서 가장 직접적으로 성립한다.
실제 다자유도 로봇에서는 방향마다 유효질량이 달라지고 축들이 서로 결합될 수 있다.
그래서 이 스칼라 공식 하나가 로봇 전체의 감쇠 특성을 모두 정해주지는 않지만, 튜닝 감각을 잡는 첫 기준으로는 매우 유용하다.

특성근으로 들어가기 전 손계산 체크를 해보자.
$m$, $b$, $k$에서 $\omega_n=\sqrt{k/m}$와 $\zeta=b/(2\sqrt{mk})$를 계산할 수 있는가?
원하는 감쇠비를 정했을 때 $b=2\zeta\sqrt{mk}$로 감쇠계수를 다시 잡을 수 있는가?
강성 $k$를 키우면 응답은 빨라질 수 있지만, 같은 $b$를 그대로 두면 감쇠비가 낮아질 수 있다는 점을 설명할 수 있는가?
이 세 질문에 답할 수 있으면, 복소근 계산은 외울 공식이 아니라 응답 모양을 분류하는 도구로 보이기 시작한다.

\subsection{특성근과 응답 형태}

지금까지는 그래프의 모양을 말로 설명했다.
이제 같은 이야기를 방정식 안에서 읽어 보자.
그 역할을 하는 것이 특성근이다.
감쇠비에 따라 응답이 달라지는 이유는 특성방정식에서 확인할 수 있다.
앞 장의 전달함수에서 pole이라고 부른 값과 여기서 특성근이라고 부르는 값은 같은 분모에서 나온다.
전달함수는 외력으로 밀었을 때의 zero-state 응답을 묻고, 자유응답은 처음 흔들어 놓고 놓았을 때를 묻지만, 둘 다 $m\ddot{x}+b\dot{x}+kx$라는 같은 왼쪽 동역학을 공유한다.
그래서 특성근을 읽는 일은 곧 pole을 읽는 일이다.
복소근 계산보다 먼저 그래프 모양을 보자.
그림~\ref{fig:damping-ratio-atlas}는 같은 표준 2차 시스템에서 감쇠비만 바꾸면 응답의 모양이 어떻게 달라지는지 보여준다.
먼저 이 그림에서 흔들림, 오버슈트, 느린 수렴의 차이를 눈으로 보고, 그다음 특성근이 왜 그런 모양을 만드는지 따라가면 된다.

\begin{figure}[!htbp]
  \centering
  \resizebox{\linewidth}{!}{%
  \begin{tikzpicture}[
      axis/.style={-{Latex[length=1.6mm]}, draw=black!65, line width=0.45pt},
      curve/.style={line width=0.95pt},
      target/.style={densely dashed, draw=black!55, line width=0.45pt},
      band/.style={densely dotted, draw=black!35, line width=0.4pt},
      paneltitle/.style={font=\scriptsize\bfseries, align=center},
      lab/.style={font=\scriptsize, align=center},
      note/.style={draw=black!35, rounded corners=2pt, fill=black!3, align=center, font=\scriptsize, inner sep=4pt}
    ]

    \begin{scope}[shift={(0,3.0)}]
      \draw[axis] (0,0) -- (3.35,0) node[right, lab] {$t$};
      \draw[axis] (0,0) -- (0,2.25);
      \draw[target] (0,1.25) -- (3.2,1.25);
      \draw[band] (0,1.31) -- (3.2,1.31);
      \draw[band] (0,1.19) -- (3.2,1.19);
      \node[paneltitle] at (1.6,2.45) {$\zeta=0$\\무감쇠};
    \end{scope}

    \begin{scope}[shift={(4.0,3.0)}]
      \draw[axis] (0,0) -- (3.35,0) node[right, lab] {$t$};
      \draw[axis] (0,0) -- (0,2.25);
      \draw[target] (0,1.25) -- (3.2,1.25);
      \draw[band] (0,1.31) -- (3.2,1.31);
      \draw[band] (0,1.19) -- (3.2,1.19);
      \node[paneltitle] at (1.6,2.45) {$\zeta=0.2$\\부족감쇠};
    \end{scope}

    \begin{scope}[shift={(8.0,3.0)}]
      \draw[axis] (0,0) -- (3.35,0) node[right, lab] {$t$};
      \draw[axis] (0,0) -- (0,2.25);
      \draw[target] (0,1.25) -- (3.2,1.25);
      \draw[band] (0,1.31) -- (3.2,1.31);
      \draw[band] (0,1.19) -- (3.2,1.19);
      \node[paneltitle] at (1.6,2.45) {$\zeta=0.7$\\부족감쇠};
    \end{scope}

    \begin{scope}[shift={(2.0,0)}]
      \draw[axis] (0,0) -- (3.35,0) node[right, lab] {$t$};
      \draw[axis] (0,0) -- (0,2.25);
      \draw[target] (0,1.25) -- (3.2,1.25);
      \draw[band] (0,1.31) -- (3.2,1.31);
      \draw[band] (0,1.19) -- (3.2,1.19);
      \node[paneltitle] at (1.6,2.45) {$\zeta=1$\\임계감쇠};
    \end{scope}

    \begin{scope}[shift={(6.0,0)}]
      \draw[axis] (0,0) -- (3.35,0) node[right, lab] {$t$};
      \draw[axis] (0,0) -- (0,2.25);
      \draw[target] (0,1.25) -- (3.2,1.25);
      \draw[band] (0,1.31) -- (3.2,1.31);
      \draw[band] (0,1.19) -- (3.2,1.19);
      \node[paneltitle] at (1.6,2.45) {$\zeta=1.5$\\과감쇠 예};
    \end{scope}

    \begin{scope}[shift={(0,3.0)}]
      \draw[curve, black!55] plot[smooth] coordinates {
        (0,0) (0.35,0.45) (0.7,1.25) (1.05,1.95) (1.4,1.75)
        (1.75,0.95) (2.1,0.25) (2.45,0.48) (2.8,1.28) (3.15,1.92)
      };
      \node[lab, text=black!60] at (1.62,0.58) {에너지 소산 없음\\정착하지 않음};
    \end{scope}

    \begin{scope}[shift={(4.0,3.0)}]
      \draw[curve, red!70!black] plot[smooth] coordinates {
        (0,0) (0.35,0.32) (0.7,0.95) (1.05,1.62) (1.35,1.86)
        (1.7,1.62) (2.05,1.2) (2.45,1.05) (2.8,1.18) (3.15,1.28)
      };
      \draw[-{Latex[length=1.6mm]}, red!70!black, line width=0.7pt] (1.35,1.25) -- (1.35,1.86);
      \node[lab, text=red!70!black] at (2.27,1.75) {큰 오버슈트};
    \end{scope}

    \begin{scope}[shift={(8.0,3.0)}]
      \draw[curve, blue!70!black] plot[smooth] coordinates {
        (0,0) (0.35,0.22) (0.75,0.67) (1.15,1.05) (1.55,1.25)
        (1.95,1.31) (2.35,1.28) (2.75,1.25) (3.15,1.25)
      };
      \node[lab, text=blue!70!black] at (2.2,1.72) {작은 오버슈트\\빠른 정착 경향};
    \end{scope}

    \begin{scope}[shift={(2.0,0)}]
      \draw[curve, green!45!black] plot[smooth] coordinates {
        (0,0) (0.35,0.12) (0.75,0.38) (1.15,0.68) (1.55,0.93)
        (1.95,1.1) (2.35,1.19) (2.75,1.23) (3.15,1.25)
      };
      \node[lab, text=green!45!black] at (2.2,0.65) {비진동 응답의\\경계};
    \end{scope}

    \begin{scope}[shift={(6.0,0)}]
      \draw[curve, orange!80!black] plot[smooth] coordinates {
        (0,0) (0.45,0.08) (0.9,0.22) (1.35,0.4) (1.8,0.58)
        (2.25,0.76) (2.7,0.92) (3.15,1.04)
      };
      \node[lab, text=orange!80!black] at (2.15,0.35) {진동은 없지만\\느린 수렴};
    \end{scope}

    \node[lab, anchor=west] at (11.35,4.25) {목표값};
    \draw[target] (10.85,4.25) -- (11.25,4.25);
    \node[lab, anchor=west, text=black!55] at (11.35,3.9) {$\pm2\%$ 정착 범위};
    \draw[band] (10.85,3.9) -- (11.25,3.9);

    \node[note, align=left] at (10.2,1.25) {표준 2차 모델의\\정규화 단위계단 응답\\$m,k$ 고정, $b$ 변화};
  \end{tikzpicture}%
  }
  \caption{감쇠비 $\zeta$가 달라질 때 표준 2차 질량-스프링-댐퍼 모델의 정규화된 단위계단 응답이 어떻게 달라지는지 보여주는 응답 아틀라스이다. 여기서는 $m$과 $k$를 고정하고 $b$를 바꾸어 $\zeta=b/(2\sqrt{mk})$가 달라지는 상황을 그렸다. 낮은 감쇠비는 오버슈트와 진동을 만들고, 임계감쇠는 비진동 응답의 경계이며, 과감쇠 응답은 이 예에서 더 느리게 수렴할 수 있다.}
  \label{fig:damping-ratio-atlas}
\end{figure}


감쇠가 작으면 흔들리며 줄어들고, 감쇠가 임계값에 가까우면 흔들리지 않고 빠르게 돌아오며, 감쇠가 너무 크면 느리게 돌아온다.
특성근은 이 세 가지 모양이 왜 생기는지를 수학적으로 보여주는 도구이다.
여기서 $x(t)=e^{st}$를 넣는 이유는 해의 모양을 외우려는 것이 아니라, 응답이 줄어드는지, 흔들리는지, 발산하는지를 한 번에 읽기 위해서이다.
이때 $e^{st}$의 $e$는 오차가 아니라 자연상수, 즉 지수함수의 밑이다.
뒤의 임피던스 제어 장에서는 $e=x-x_d$처럼 오차를 $e$로 쓰므로 문맥을 구분해야 한다.
$s$는 실수일 수도 있고 복소수일 수도 있다.
예를 들어 $s=-2$이면 그냥 줄어드는 지수함수이고, $s=-2+3\jj$이면 줄어들면서 흔들리는 함수가 된다.
미분해도 모양이 크게 바뀌지 않는 함수가 지수함수이기 때문이다.
$e^{st}$를 한 번 미분하면 $se^{st}$, 두 번 미분하면 $s^2e^{st}$가 되므로, 미분방정식이 $s$에 대한 대수방정식으로 바뀐다.
즉 복잡한 시간 함수의 모양을 직접 찾기 전에, 가능한 응답의 성장, 감소, 진동 성질을 $s$라는 숫자로 분류하는 것이다.
자유응답 $f(t)=0$을 가정하고 $x(t)=e^{st}$를 대입하면
\begin{equation}
  ms^2+bs+k=0
\end{equation}
을 얻는다.
이 방정식의 해 $s$를 특성근(characteristic root)이라고 한다.
특성근에서 먼저 읽어야 할 것은 응답 곡선의 모양이다.
세 가지만 붙잡으면 충분하다.
\begin{itemize}
  \item 특성근이 서로 다른 실수이면 흔들림 없이 지수함수들의 합으로 돌아온다.
  \item 특성근이 복소수이면 사인과 코사인 성분이 생기므로 흔들린다.
  \item 특성근의 실수부가 음수이면 시간이 지날수록 줄어들고, 양수이면 발산한다.
\end{itemize}
먼저 이차방정식 공식으로 쓰면
\begin{equation}
  s_{1,2}
  =
  \frac{-b\pm\sqrt{b^2-4mk}}{2m}
\end{equation}
이다.
여기에 표준 파라미터 $\omega_n=\sqrt{k/m}$, $\zeta=b/(2\sqrt{mk})$를 넣으면 특성근이 훨씬 읽기 좋은 모양으로 정리된다.
이 대입은 논문에서 보통 한 줄로 건너뛰지만, 두 조각으로 나누면 손으로 따라갈 수 있다.
먼저 $\zeta=b/(2\sqrt{mk})$를 $b$에 대해 뒤집으면 $b=2\zeta\sqrt{mk}$이다.
첫 조각인 분수 앞부분 $-b/(2m)$은
\[
  -\frac{b}{2m}
  =
  -\frac{2\zeta\sqrt{mk}}{2m}
  =
  -\zeta\frac{\sqrt{mk}}{m}
  =
  -\zeta\sqrt{\frac{mk}{m^2}}
  =
  -\zeta\sqrt{\frac{k}{m}}
  =
  -\zeta\omega_n
\vmark{charroot-substitution-decay-piece}%
\]
이 된다.
둘째 조각인 근호 안은 $b^2=4\zeta^2mk$이므로
\[
  b^2-4mk
  =
  4\zeta^2mk-4mk
  =
  4mk(\zeta^2-1)
\vmark{charroot-substitution-radical-piece}%
\]
이고, 제곱근을 취하면 $\sqrt{b^2-4mk}=2\sqrt{mk}\sqrt{\zeta^2-1}$, 이를 분모 $2m$으로 나누면 첫 조각과 같은 계산으로 $\sqrt{k/m}\,\sqrt{\zeta^2-1}=\omega_n\sqrt{\zeta^2-1}$이 된다.
두 조각을 합치면 특성근은
\begin{equation}
  s_{1,2}
  =
  -\zeta\omega_n
  \pm
  \omega_n\sqrt{\zeta^2-1}
\end{equation}
이다.
숫자로 검산해 보자.
$m=1~\mathrm{kg}$, $b=2~\mathrm{N\,s/m}$, $k=4~\mathrm{N/m}$이면 $\omega_n=\sqrt{4/1}=2~\mathrm{rad/s}$, $\zeta=2/(2\sqrt{4})=0.5$이다.
이차방정식 공식으로 직접 풀면 $s_{1,2}=(-2\pm\sqrt{4-16})/2=-1\pm\sqrt{-3}$이고, 표준형으로 계산하면 $-\zeta\omega_n=-1$, $\omega_n\sqrt{\zeta^2-1}=2\sqrt{-0.75}=\sqrt{-3}$이므로 두 계산이 정확히 일치한다.
\vmark{charroot-numeric-check}%
이 식에서 앞의 $-\zeta\omega_n$은 응답이 공통적으로 줄어드는 속도와 관련되고, 제곱근이 붙은 뒤쪽 항은 근이 실수인지 복소수인지, 즉 흔들리는지 아닌지를 가른다.
따라서 특성근은 그래프의 감소 속도와 진동 여부를 담은 응답 모양의 지문처럼 읽으면 된다.

부족감쇠 조건 $0<\zeta<1$에서는 루트 안의 값이 음수가 되므로 근이 복소수가 된다.
$\jj$는 제곱하면 $-1$이 되는 허수 단위이다.
음수의 제곱근은 이 $\jj$를 밖으로 꺼내 쓸 수 있다.
일반적으로 양수 $y$에 대해 $\sqrt{-y}=\sqrt{(-1)\,y}=\sqrt{-1}\,\sqrt{y}=\jj\sqrt{y}$로 읽으면 된다.
$\zeta^2-1<0$이므로 $\zeta^2-1=(-1)(1-\zeta^2)$로 쪼개면
\[
  \sqrt{\zeta^2-1}
  =
  \sqrt{(-1)(1-\zeta^2)}
  =
  \jj\sqrt{1-\zeta^2}
\vmark{charroot-imaginary-unit-factorization}%
\]
이고, 여기서 $1-\zeta^2>0$이라 마지막 제곱근은 보통의 양수 제곱근이다.
따라서 특성근의 뒤쪽 항은 $\omega_n\sqrt{\zeta^2-1}=\jj\,\omega_n\sqrt{1-\zeta^2}$가 되고, 이 허수부의 크기에 이름을 붙인 것이 감쇠 진동수이다.
\begin{equation}
  s_{1,2}
  =
  -\zeta\omega_n
  \pm
  \jj \omega_d
\end{equation}
여기서
\begin{equation}
  \omega_d=\omega_n\sqrt{1-\zeta^2}
\end{equation}
는 감쇠 진동수(damped natural frequency)이다.
특성근에 허수부가 있다는 것은 시간응답에 사인과 코사인 형태의 흔들림이 생긴다는 뜻이다.
부족감쇠 자유응답은
\begin{equation}
  x(t)=e^{-\zeta\omega_n t}
  \left(A\cos\omega_d t+B\sin\omega_d t\right)
\end{equation}
형태가 된다.
여기서 $A$와 $B$는 처음 위치와 처음 속도에 의해 정해지는 상수이다.
같은 시스템이라도 처음에 얼마나 당겼는지, 손을 놓을 때 속도가 있었는지에 따라 진동의 시작 모양은 달라진다.

이 식에서 응답의 모양이 바로 보인다.
앞의
\begin{equation}
  e^{-\zeta\omega_n t}
\end{equation}
는 진폭을 감싸는 포락선(envelope)이다.
포락선은 진동의 꼭대기들을 이어서 감싸는 선이라고 생각하면 된다.
시간이 지날수록 이 값이 작아지므로 진동의 크기가 줄어든다.
뒤의
\begin{equation}
  A\cos\omega_d t+B\sin\omega_d t
\end{equation}
는 실제로 왔다 갔다 흔들리는 진동 부분이다.
부족감쇠 응답은 줄어드는 포락선 안에서 흔들리는 운동이다.

감쇠비가 정확히 $\zeta=1$이면 두 특성근이 모두
\begin{equation}
  s=-\omega_n
\end{equation}
으로 겹치는 중근이 된다.
이때 자유응답은
\begin{equation}
  x(t)=(A+Bt)e^{-\omega_n t}
\end{equation}
형태이며, 복소근에 의한 반복 진동 없이 평형점으로 돌아오는 경계 응답이다.
$\zeta>1$이면 두 특성근은 서로 다른 음의 실수가 된다.
이 경우도 흔들리지는 않지만, 두 지수응답 중 더 천천히 줄어드는 항이 전체 정착 속도를 지배하므로 임계감쇠보다 느리게 보일 수 있다.

부족감쇠 범위 $0<\zeta<1$에서는 감쇠비가 커질수록 두 가지 일이 동시에 일어난다.
첫째, $e^{-\zeta\omega_n t}$가 더 빨리 줄어들기 때문에 진폭이 빨리 죽는다.
둘째, $\omega_d=\omega_n\sqrt{1-\zeta^2}$가 작아지므로 실제 진동 주파수는 조금 낮아진다.
$\zeta$가 1에 가까워지면 더 이상 진동하지 않는 경계에 도달한다.
반면 $\zeta\ge1$에서는 $\omega_d$를 실제 진동 주파수로 해석하지 않는다.
이때 응답은 복소근이 아니라 실수근들이 만드는 지수응답으로 돌아오며, 감쇠가 너무 커지면 오히려 가장 느린 실수근이 전체 수렴을 지배해 움직임이 둔해질 수 있다.

\subsection{오버슈트와 정착시간의 의미}

제어에서 감쇠비와 고유진동수를 중요하게 보는 이유는 시간응답 특성과 직접 연결되기 때문이다.
대표적인 지표는 상승시간(rise time), 오버슈트(overshoot), 정착시간(settling time), 시정수(time constant), 정상상태 오차(steady-state error)이다.
여기서도 먼저 그래프를 떠올리는 것이 좋다.
목표값이 수평선으로 그려져 있고, 실제 응답 곡선이 그 선을 넘어갔다가 다시 내려오면 오버슈트가 생긴 것이다.
응답 곡선이 목표 주변의 작은 띠 안에 들어간 뒤 더 이상 크게 벗어나지 않으면 정착했다고 말한다.
충분히 시간이 지난 뒤 목표값과 최종 응답 사이에 남아 있는 차이는 정상상태 오차이다.
조금 더 물리적으로 말하면, 과도응답은 시스템이 최종 평형점으로 가는 동안의 이야기이고 정상상태응답은 속도와 가속도가 0이 된 뒤의 이야기이다.
일정한 힘 $F_0$가 질량-스프링-댐퍼에 계속 걸리면 충분히 오래 지난 뒤에는
\begin{equation}
  \dot{x}=0,\qquad \ddot{x}=0
\end{equation}
이므로 스프링 항만 남아 $kx_{\mathrm{ss}}=F_0$가 된다.
따라서
\begin{equation}
  x_{\mathrm{ss}}=\frac{F_0}{k}
\end{equation}
이다.
오버슈트와 정착시간은 이 최종값에 도달하는 과정에서 얼마나 지나치고 얼마나 오래 흔들리는지를 말해준다.

\begin{table}[!htbp]
  \centering
  \caption{시간응답 그래프에서 자주 읽는 지표}
  \small
  \setlength{\tabcolsep}{4pt}
  \begin{tabularx}{\linewidth}{@{}>{\raggedright\arraybackslash}p{0.18\linewidth}>{\raggedright\arraybackslash}p{0.31\linewidth}>{\raggedright\arraybackslash}X@{}}
    \toprule
    지표 & 그래프에서 읽는 법 & 주로 말해주는 것 \\
    \midrule
    상승시간 & 응답이 최종값의 10\%에서 90\%까지 올라가는 데 걸리는 시간. 문헌에 따라 0--100\% 첫 도달을 쓰기도 하므로 기준을 밝혀야 한다. & 목표 근처까지 얼마나 빨리 접근하는가 \\
    최대 오버슈트 & 최대 피크값에서 최종값을 뺀 뒤, 보통 최종값으로 나누어 비율로 표시한다. & 목표를 얼마나 지나치는가 \\
    정착시간 & 응답이 최종값 주변의 $\pm2\%$ 또는 $\pm5\%$ 띠 안에 들어간 뒤 계속 머무는 첫 시각. & 흔들림이 실질적으로 사라지는 데 걸리는 시간 \\
    시정수 & 1차 응답에서는 최종값의 약 63.2\%에 도달하는 시간이다. 지수 감쇠 $e^{-t/\tau}$에서는 포락선이 약 36.8\%로 줄어드는 시간이다. & 지수적으로 빨라지거나 줄어드는 기본 시간 척도 \\
    정상상태 오차 & 충분히 시간이 지난 뒤 목표값과 실제 최종값의 차이. & 오래 기다려도 남는 위치 또는 출력 오차 \\
    \bottomrule
  \end{tabularx}
\end{table}

오버슈트는 목표값을 지나쳐 얼마나 더 올라갔는지를 나타낸다.
먼저 그래프에서는 최대 피크값과 최종값의 차이로 읽는다.
그다음 표준 2차 시스템이라는 조건이 맞을 때, 이 값을 감쇠비 하나로 빠르게 예측하는 대표 공식이 있다.
아래 식은 모든 시스템에 항상 정확히 맞는 만능 공식이 아니라, 부족감쇠 표준 2차 시스템의 단위계단 응답에서 자주 쓰는 대표 공식이다.
여기서 단위계단 응답은 입력이나 목표를 0에서 1로 갑자기 바꾼 뒤 그대로 유지했을 때의 응답이다.
영 초기조건은 응답이 시작하기 전 출력과 속도 같은 초기 상태가 0이라고 두는 조건이다.
더 구체적으로는
\begin{equation}
  G(s)=\frac{\omega_n^2}{s^2+2\zeta\omega_n s+\omega_n^2}
\end{equation}
처럼 쓸 수 있는 표준 2차 시스템, 영 초기조건, 양의 단위계단 입력, $0<\zeta<1$, 추가적인 zero나 포화가 없는 경우를 생각한다.
그때 최대 오버슈트 비율은
\begin{equation}
  M_p
  =
  \exp\left(
  -\frac{\zeta\pi}{\sqrt{1-\zeta^2}}
  \right)
\end{equation}
로 나타낼 수 있다.
이 식의 $M_p$는 $0.046$ 같은 비율이다.
퍼센트로 말할 때는 $100M_p\%$로 바꾼다.
그래프에서 읽을 때는 최대 피크값 $y_{\mathrm{peak}}$와 최종값 $y_{\infty}$를 잡고,
\begin{equation}
  M_p=\frac{y_{\mathrm{peak}}-y_{\infty}}{y_{\infty}}
\end{equation}
처럼 계산한다고 생각하면 된다.
여기서 $y$는 그래프에서 관찰하는 출력이다.
이 장에서는 보통 위치 $x$를 떠올리면 되고, $y_{\infty}$는 과도응답이 사라진 뒤 충분히 오래 기다렸을 때 남는 최종값이다.
양의 최종값을 갖는 표준 단위계단 응답에서는 위 정의가 자연스럽다.
일반 출력에서는 보통 $|y_{\infty}|$로 정규화하거나, 목표값이 0에 가깝거나 부호가 바뀌는 경우에는 퍼센트 오버슈트의 정의를 따로 정해야 한다.
감쇠비 $\zeta$가 작으면 오버슈트가 크고, $\zeta$가 커지면 오버슈트가 줄어든다.
물리적으로는 감쇠가 작을수록 에너지가 많이 남아 있어서 목표점을 지나친 뒤에도 계속 움직이기 때문이다.
예를 들어 $\zeta=0.2$이면 이 공식은 약 $53\%$의 오버슈트를 예측한다.
목표가 $1$이라면 순간적으로 약 $1.53$까지 올라갈 수 있다는 뜻이다.
반면 $\zeta=0.7$이면 오버슈트가 약 $4.6\%$ 수준으로 줄어든다.
이 값이 어떻게 나오는지 한 번만 대입 과정을 보자.
$\zeta=0.7$이면 $\sqrt{1-\zeta^2}=\sqrt{1-0.49}=\sqrt{0.51}\approx0.714$이고, 지수 안은 $-0.7\pi/0.714\approx-3.08$이므로 $M_p=e^{-3.08}\approx0.046$, 즉 약 $4.6\%$이다.
\vmark{overshoot-zeta07-substitution}%
이 숫자에서 읽어야 할 것은 감쇠비가 그래프의 출렁임과 직접 연결된다는 점이다.

정착시간은 응답이 목표값 근처의 작은 범위 안에 들어간 뒤 계속 머무는 데 걸리는 시간이다.
다음 식은 주로 $0<\zeta<1$인 부족감쇠 표준 2차계에서 빠른 설계 감각을 잡기 위한 근사이다.
목표값의 $\pm2\%$ 띠를 기준으로 자주 쓰는 대략적인 정착시간은
\begin{equation}
  T_s \approx \frac{4}{\zeta\omega_n}
\end{equation}
으로 많이 사용된다.
이 식은 $\zeta\omega_n$이 클수록 진폭 포락선이 빠르게 줄어든다는 뜻이다.
왜 4가 나오는지도 감으로 볼 수 있다.
엄밀한 표준 2차 단위계단 오차 포락선에는 $1/\sqrt{1-\zeta^2}$ 계수가 함께 붙지만, 중간 감쇠비 범위의 설계 감각에서는 먼저 $e^{-\zeta\omega_n t}$만 남겨 단순화해 볼 수 있다.
부족감쇠 응답의 진폭 포락선은 대략 $e^{-\zeta\omega_n t}$처럼 줄어든다.
이 포락선이 2\% 수준, 즉 $0.02=1/50$ 정도로 줄어드는 시간을 잡으면
\begin{equation}
  e^{-\zeta\omega_n t}\approx0.02
\end{equation}
이고, 양변에 로그를 취하면 $\zeta\omega_n t\approx\ln 50\approx3.9$가 된다.
$e^{3.9}$가 약 50이기 때문에, 그만큼 시간이 지나면 포락선이 약 50분의 1로 줄었다고 읽을 수 있다.
그래서 $4/(\zeta\omega_n)$라는 기억하기 쉬운 근사가 나온다.
즉, 고유진동수가 크고 감쇠비도 적절히 크면 빠르게 안정된다.
하지만 감쇠를 무작정 크게 하면 과감쇠가 되어 오히려 느려질 수 있다.
또 이 근사는 2\% 기준의 표준 부족감쇠 2차계에서 주로 쓰는 포락선 근사이다.
실무적으로는 대략 $0.4\lesssim\zeta\lesssim0.8$ 정도의 설계 감각에서 자주 쓰이며, $\zeta$가 1에 가까워지거나 매우 작아지면 실제 피크 시점과 위상 때문에 차이가 커질 수 있다.
5\% 기준을 쓰면 상수가 대략 3에 가까워지고, $\zeta\ge1$인 과감쇠 시스템이나 고차 시스템, 포화와 마찰이 강한 실제 로봇에는 그대로 들어맞지 않을 수 있다.
그래서 제어 설계에서는 빠른 응답과 작은 오버슈트 사이의 균형을 잡아야 한다.

앞에서 본 숫자를 이어서 써 보자.
$\omega_n=10~\mathrm{rad/s}$이고 $\zeta=0.7$인 표준 2차 시스템이라면, 오버슈트 공식은 약 $4.6\%$를 예측한다.
목표 위치가 $1~\mathrm{m}$라면 최대 피크가 대략 $1.046~\mathrm{m}$ 근처라는 뜻이다.
2\% 정착 기준이라면 $0.98~\mathrm{m}$에서 $1.02~\mathrm{m}$ 사이에 들어간 뒤 계속 머물러야 정착했다고 부른다.
또 2\% 정착시간 근사는
\begin{equation}
  T_s \approx \frac{4}{0.7\cdot10}\approx0.57~\mathrm{s}
\end{equation}
이다.
이 숫자들은 정확한 시뮬레이터 결과를 대신하는 값이 아니라, 그래프를 보기 전에 응답이 어느 정도 빠르고 얼마나 출렁일지 예상하는 기준선이다.

제어를 배울 때는 이 분류를 꼭 한 번 지나가게 된다.
학습자는 같은 모델에서 $m$, $b$, $k$를 바꾸면서 오버슈트, 정착시간, 정상상태 오차, 제어입력 크기가 어떻게 달라지는지 직접 확인할 수 있다.
Lab01은 MuJoCo 슬라이드 관절 모델을 이용해 이 질량-스프링-댐퍼 응답을 재현하도록 설계되어 있다.
따라서 이론식은 그래프를 보기 전의 예측표이다.
예를 들어 $k$를 키우기 전에는 진동 주기가 짧아질지, $b$를 키우기 전에는 오버슈트가 줄어들지 먼저 예상하고, Lab01 플롯에서 그 변화를 확인한다.
로봇 임피던스 튜닝에서도 같은 방식으로 생각하면 된다.
강성을 올리면 목표로 돌아가려는 힘은 커지지만, 감쇠를 함께 조정하지 않으면 끝단이 목표를 지나쳐 흔들릴 수 있다.
그래서 오버슈트와 정착시간은 단순한 그래프 지표가 아니라, 로봇이 빠르게 느껴지는지, 부드럽게 느껴지는지, 불안하게 떨리는지를 판단하는 튜닝 언어가 된다.

같은 운동방정식이라도 출력과 비율을 어떻게 잡느냐에 따라 이름이 달라진다.
위치를 출력으로 보면 순응성 $X(s)/F(s)$이고, 영 초기조건의 전달함수 문맥에서 속도를 출력으로 보면 모빌리티 $V_x(s)/F(s)$이다.
반대로 주어진 속도 운동을 만들기 위해 필요한 힘을 묻으면 기계적 임피던스 $F(s)/V_x(s)$가 된다.

\subsection{기계적 임피던스 관점}

질량-스프링-댐퍼 시스템의 기계적 임피던스는
\begin{align}
  F(s)
  &=
  \left(ms^2+bs+k\right)X(s), \\
  V_x(s)
  &=
  sX(s)
\end{align}
에서 바로 얻을 수 있다.
첫 번째 식은 운동방정식 $m\ddot{x}+b\dot{x}+kx=f$를 영 초기조건에서 라플라스 변환한 것이고, 두 번째 식은 속도가 위치의 시간 미분이라는 사실을 같은 영 초기조건의 전달함수 문맥으로 옮긴 것이다.
초기 위치가 0이 아니면 엄밀히는 $V_x(s)=sX(s)-x(0)$가 되므로, 이미 당겨 놓고 시작하는 자유응답은 따로 구분해야 한다.
따라서 힘을 속도로 나누면
\begin{equation}
  Z_m(s)=\frac{F(s)}{sX(s)}
  =
  ms+b+\frac{k}{s}
\end{equation}
이다.
앞 장에서는 이 식을 전기 임피던스와 나란히 놓기 위한 예고편으로 보았다.
이제는 실제 기계 요소의 관점에서 다시 읽을 수 있다.
특히 스프링 항이 $k/s$로 보이는 이유를 한 번 더 붙잡아 두자.
시간영역에서 외부에서 필요한 스프링 대응 힘은 $+kx$이고, 물체에 작용하는 실제 복원력은 $-kx$이다.
임피던스 식의 $F$는 주어진 운동을 만들기 위해 외부에서 넣어야 하는 힘 쪽으로 읽는다.
그래서 스프링은 외부 힘과 위치의 관계로 보면 $f=kx$처럼 보인다.
하지만 임피던스는 힘을 위치가 아니라 속도로 나눈 값이다.
같은 변위를 아주 천천히 만들면 속도가 작으므로 힘/속도 비가 크게 보이고, 그래서 낮은 주파수에서 스프링 항이 크게 나타난다.
진폭만 놓고 보면 더 직접적이다.
$x(t)=X_0\cos(\omega t)$로 움직일 때 속도 진폭은 $\omega X_0$이고, 스프링 힘의 진폭은 $kX_0$이다.
따라서 스프링만의 힘/속도 진폭비는
\begin{equation}
  \frac{kX_0}{\omega X_0}=\frac{k}{\omega}
\end{equation}
가 된다.
이것이 주파수 영역에서 스프링 항이 $k/\omega$처럼 보이는 이유이다.
주파수 응답으로 쓰면
\begin{equation}
  Z_m(\jj\omega)
  =
  b+\jj\left(\omega m-\frac{k}{\omega}\right).
\end{equation}
이 식을 읽을 때는 힘을 속도로 나눈다는 말을 손으로 미는 느낌으로 바꾸어 보면 좋다.
천천히 밀어 위치만 바꾸려 하면 스프링이 주로 버틴다.
그래서 낮은 주파수에서는 $k/\omega$가 크게 보인다.
반대로 빠르게 흔들면 위치가 조금만 움직여도 가속도가 커지고, 질량을 움직이는 힘이 커진다.
그래서 높은 주파수에서는 $\omega m$이 커진다.
댐퍼는 어느 경우에도 속도에 비례해 에너지를 빼앗으므로 실수부 $b$로 남는다.
허수부의 질량항과 스프링항이 서로 다른 부호를 갖는 것은 한쪽이 좋은 효과이고 다른 쪽이 나쁜 효과라는 뜻이 아니라, 저장했다가 돌려주는 에너지의 위상 방향이 서로 반대라는 뜻이다.
\begin{itemize}
  \item 감쇠항 $b$는 주파수와 무관한 실수 성분이다. 속도에 비례해 에너지를 소산한다.
  \item 질량항 $\omega m$은 주파수가 높을수록 커진다. 빠르게 흔들수록 큰 관성력이 필요하다.
  \item 스프링항 $k/\omega$는 낮은 주파수에서 커진다. 천천히 위치를 밀어 바꾸려 할수록 복원력이 지배적이다.
\end{itemize}
질량항도 같은 방식으로 볼 수 있다.
$x(t)=X_0\cos(\omega t)$이면 가속도 진폭은 $\omega^2X_0$이고, 관성에 필요한 힘의 진폭은 $m\omega^2X_0$이다.
이를 속도 진폭 $\omega X_0$로 나누면 $m\omega$가 된다.
그래서 빠른 진동일수록 질량항이 커진다.
댐퍼는 더 단순하다.
댐퍼 힘의 진폭은 $b\omega X$이고 속도 진폭도 $\omega X$이므로, 힘/속도 비는 항상 $b$이다.
따라서 댐퍼는 주파수에 따라 저장 성질이 바뀌기보다, 속도에 비례한 소산항으로 남는다.
그림~\ref{fig:mechanical-impedance-frequency}는 이 세 항을 주파수에 따라 나누어 보는 정성적 지도이다.
숫자를 정확히 읽기 위한 그래프라기보다, 느린 운동에서는 스프링 효과가, 빠른 운동에서는 질량 효과가 왜 눈에 띄는지 이해하기 위한 그림이다.
다만 이 그림은 전체 임피던스 크기 자체가 아니라 개별 항의 크기를 비교하는 그림이다.
실제 전체 크기는 $|Z_m(\jj\omega)|=\sqrt{b^2+(\omega m-k/\omega)^2}$처럼 스프링 항과 질량 항의 부호까지 함께 고려해야 한다.

\begin{figure}[!htbp]
  \centering
  \resizebox{0.92\linewidth}{!}{%
  \begin{tikzpicture}[
      axis/.style={-{Latex[length=2.2mm]}, draw=black!70, thick},
      curve/.style={line width=1.0pt},
      region/.style={font=\scriptsize, align=center, text=black!70},
      callout/.style={draw=black!35, rounded corners=2pt, fill=black!3, align=center, font=\scriptsize, inner sep=4pt}
    ]
    \draw[axis] (0,0) -- (9.0,0) node[right, font=\small] {각주파수 $\omega$};
    \draw[axis] (0,0) -- (0,4.9) node[above, font=\small] {개별 항 크기};
    \node[region] at (1.45,-0.35) {느린 변화\\저주파};
    \node[region] at (4.5,-0.35) {중간 영역};
    \node[region] at (7.6,-0.35) {빠른 변화\\고주파};

    \fill[blue!5] (0.2,0.05) rectangle (2.8,4.55);
    \fill[green!5] (2.8,0.05) rectangle (6.0,4.55);
    \fill[red!5] (6.0,0.05) rectangle (8.8,4.55);

    \draw[curve, blue!70!black] plot[smooth] coordinates {
      (0.55,4.35) (1.05,3.55) (1.75,2.72) (2.55,2.02)
      (3.55,1.35) (4.75,0.9) (6.0,0.65) (8.45,0.42)
    };
    \draw[curve, green!55!black, dashed] (0.55,2.05) -- (8.45,2.05);
    \draw[curve, red!70!black] plot[smooth] coordinates {
      (0.55,0.42) (1.45,0.62) (2.45,0.95) (3.6,1.45)
      (4.9,2.12) (6.25,2.95) (7.45,3.65) (8.45,4.25)
    };
    \draw[densely dashed, black!45] (4.45,0.1) -- (4.45,4.35);
    \node[region] at (4.45,4.65) {$\omega_n$ 근처};

    \node[callout, text=blue!70!black] at (1.55,4.55) {스프링 항\\$k/\omega$};
    \node[callout, text=green!45!black] at (4.55,2.48) {감쇠 항\\$b$};
    \node[callout, text=red!70!black] at (7.25,4.45) {질량 항\\$\omega m$};
    \node[callout] at (4.35,0.58) {스프링 항과 질량 항은\\허수부에서 반대 부호\\전체 $|Z_m|$는 단순 합이 아님};
  \end{tikzpicture}%
  }
  \caption{질량-스프링-댐퍼 시스템의 기계적 임피던스는 하나의 강성값으로 고정되지 않는다. 낮은 주파수에서는 스프링 항 $k/\omega$가 크게 보이고, 높은 주파수에서는 질량 항 $\omega m$이 크게 보이며, 감쇠 항 $b$는 에너지 소산을 담당한다. 이 그림은 전체 임피던스 크기 $|Z_m(\jj\omega)|$가 아니라 개별 항 크기를 보여주는 정성적 지도이다. 실제 전체 크기는 $b$와 반응성 성분 $\omega m-k/\omega$가 함께 결정한다.}
  \label{fig:mechanical-impedance-frequency}
\end{figure}


따라서 임피던스는 단순히 얼마나 딱딱한가만을 뜻하지 않는다.
주파수에 따라 스프링처럼 보일 수도 있고, 질량처럼 보일 수도 있으며, 감쇠처럼 에너지를 소산할 수도 있다.
로봇 임피던스 제어에서 강성, 감쇠, 관성을 따로 설계하는 이유가 바로 여기에 있다.

로봇 장으로 넘어갈 때는 먼저 1축 또는 모드별로 분리해 볼 수 있고, 방향별 유효질량 $m_{\mathrm{eff}}$가 정해졌다는 근사에서 아래 문장을 읽자.
\begin{itemize}
  \item 강성 $k$는 접촉 중 위치 오차와 힘의 관계를 정한다.
  \item 감쇠비 $\zeta$는 목표점 주변에서 흔들릴지, 얼마나 부드럽게 멈출지를 정한다.
  \item 고유진동수 $\omega_n$은 시스템이 얼마나 빠르게 반응하려는지를 보여주지만, 실제 로봇에서는 구동기 한계와 안정성 때문에 감쇠와 함께 봐야 한다.
\end{itemize}
이제 이 세 요소를 실제 부품이 아니라 제어기가 만들어내는 느낌으로 옮겨 보자.
로봇 손끝에 진짜 스프링과 댐퍼를 붙이지 않아도, 위치 오차와 속도를 읽어 같은 효과를 내는 힘이나 토크를 계산할 수 있다.
로봇 손끝의 한 방향만 떼어 보면 실제 또는 조작공간 유효질량 $m_{\mathrm{eff}}$ 위에 가상 강성 $k$와 감쇠 $b$를 얹은 1자유도 모델로 근사할 수 있다.
완전한 토크 제어와 동역학 보상이 있을 때는 관성까지 어느 정도 형상화할 수 있지만, 단순 임피던스 또는 PD 튜닝에서는 $m_{\mathrm{eff}}$를 주로 로봇 자세와 방향이 정하는 값으로 보고 $k$와 $b$를 조정하는 편이 안전하다.
그렇게 손끝에 원하는 힘-운동 관계를 만들려는 방법이 임피던스 제어이다.
